% Options for packages loaded elsewhere
\PassOptionsToPackage{unicode}{hyperref}
\PassOptionsToPackage{hyphens}{url}
\documentclass[
  11pt,
]{article}
\usepackage{xcolor}
\usepackage[left=3cm,right=3cm,top=2.5cm,bottom=2.5cm]{geometry}
\usepackage{amsmath,amssymb}
\setcounter{secnumdepth}{5}
\usepackage{iftex}
\ifPDFTeX
  \usepackage[T1]{fontenc}
  \usepackage[utf8]{inputenc}
  \usepackage{textcomp} % provide euro and other symbols
\else % if luatex or xetex
  \usepackage{unicode-math} % this also loads fontspec
  \defaultfontfeatures{Scale=MatchLowercase}
  \defaultfontfeatures[\rmfamily]{Ligatures=TeX,Scale=1}
\fi
\usepackage{lmodern}
\ifPDFTeX\else
  % xetex/luatex font selection
\fi
% Use upquote if available, for straight quotes in verbatim environments
\IfFileExists{upquote.sty}{\usepackage{upquote}}{}
\IfFileExists{microtype.sty}{% use microtype if available
  \usepackage[]{microtype}
  \UseMicrotypeSet[protrusion]{basicmath} % disable protrusion for tt fonts
}{}
\usepackage{setspace}
\makeatletter
\@ifundefined{KOMAClassName}{% if non-KOMA class
  \IfFileExists{parskip.sty}{%
    \usepackage{parskip}
  }{% else
    \setlength{\parindent}{0pt}
    \setlength{\parskip}{6pt plus 2pt minus 1pt}}
}{% if KOMA class
  \KOMAoptions{parskip=half}}
\makeatother
\usepackage{longtable,booktabs,array}
\usepackage{caption}
\captionsetup[table]{skip=6pt}
\newcounter{none} % for unnumbered tables
\usepackage{calc} % for calculating minipage widths
% Correct order of tables after \paragraph or \subparagraph
\usepackage{etoolbox}
\makeatletter
\patchcmd\longtable{\par}{\if@noskipsec\mbox{}\fi\par}{}{}
\makeatother
\usepackage{graphicx}
\makeatletter
\newsavebox\pandoc@box
\newcommand*\pandocbounded[1]{% scales image to fit in text height/width
  \sbox\pandoc@box{#1}%
  \Gscale@div\@tempa{\textheight}{\dimexpr\ht\pandoc@box+\dp\pandoc@box\relax}%
  \Gscale@div\@tempb{\linewidth}{\wd\pandoc@box}%
  \ifdim\@tempb\p@<\@tempa\p@\let\@tempa\@tempb\fi% select the smaller of both
  \ifdim\@tempa\p@<\p@\scalebox{\@tempa}{\usebox\pandoc@box}%
  \else\usebox{\pandoc@box}%
  \fi%
}
% Set default figure placement to htbp
\def\fps@figure{htbp}
\makeatother
% definitions for citeproc citations
\NewDocumentCommand\citeproctext{}{}
\NewDocumentCommand\citeproc{mm}{%
  \begingroup\def\citeproctext{#2}\cite{#1}\endgroup}
\makeatletter
 % allow citations to break across lines
 \let\@cite@ofmt\@firstofone
 % avoid brackets around text for \cite:
 \def\@biblabel#1{}
 \def\@cite#1#2{{#1\if@tempswa , #2\fi}}
\makeatother
\newlength{\cslhangindent}
\setlength{\cslhangindent}{1.5em}
\newlength{\csllabelwidth}
\setlength{\csllabelwidth}{3em}
\newenvironment{CSLReferences}[2] % #1 hanging-indent, #2 entry-spacing
 {\begin{list}{}{%
  \setlength{\itemindent}{0pt}
  \setlength{\leftmargin}{0pt}
  \setlength{\parsep}{0pt}
  % turn on hanging indent if param 1 is 1
  \ifodd #1
   \setlength{\leftmargin}{\cslhangindent}
   \setlength{\itemindent}{-1\cslhangindent}
  \fi
  % set entry spacing
  \setlength{\itemsep}{#2\baselineskip}}}
 {\end{list}}
\usepackage{calc}
\newcommand{\CSLBlock}[1]{\hfill\break\parbox[t]{\linewidth}{\strut\ignorespaces#1\strut}}
\newcommand{\CSLLeftMargin}[1]{\parbox[t]{\csllabelwidth}{\strut#1\strut}}
\newcommand{\CSLRightInline}[1]{\parbox[t]{\linewidth - \csllabelwidth}{\strut#1\strut}}
\newcommand{\CSLIndent}[1]{\hspace{\cslhangindent}#1}
\setlength{\emergencystretch}{3em} % prevent overfull lines
\providecommand{\tightlist}{%
  \setlength{\itemsep}{0pt}\setlength{\parskip}{0pt}}
% Statistics in Medicine submission preamble.
% Shared across analysis/report/<paper>/report.Rmd files via
%   includes:
%     in_header: ../sim-preamble.tex
% in the YAML output block.
%
% The page footer (source path, render time, version) is no
% longer set here. It is supplied at render time by the vendored
% tools/stamp-render.R, which injects tools/stamp.tex.

% Continuous line numbering for editorial review. The 'mathlines'
% option extends numbering through display-math environments
% (equation, align, ...), which SIM editorial workflows expect.
\usepackage[mathlines]{lineno}
\linenumbers
\usepackage{amsmath}

% Spacing: linestretch is set in YAML (1.5 line spacing).
\usepackage{setspace}
%% tools/stamp.tex
%% Provenance footer for PDFs rendered from this project.
%% Vendored by zzcollab; this file is not project-specific. The
%% footer prints the source document and the time of rendering.
%% The git version is defined here too, and so is available to a
%% project that wants it on the page, but it is left off the
%% default footer: it already names the staged copy in share/ and
%% appears in its manifest row. All values are supplied at render
%% time by tools/stamp-render.R, which writes a small generated
%% file that redefines the macros below. The \providecommand
%% defaults keep a bare LaTeX compile working when the document is
%% built without the wrapper.
\usepackage{fancyhdr}
\usepackage{xcolor}
\providecommand{\stampsource}{\detokenize{(source not recorded)}}
\providecommand{\stamptime}{(time not recorded)}
\providecommand{\stampversion}{\detokenize{(version not recorded)}}
\setlength{\footskip}{40pt}
\newcommand{\stampfootcontent}{%
  \footnotesize\color{gray}%
  \parbox[b]{\textwidth}{%
    \stamptime\hfill\thepage\\[2pt]%
    \ttfamily\raggedright\stampsource}}
\pagestyle{fancy}
\fancyhf{}
\fancyfoot[C]{\stampfootcontent}
\renewcommand{\headrulewidth}{0pt}
\renewcommand{\footrulewidth}{0.4pt}
%% The title page uses the `plain` page style; redefine it so the
%% stamp appears there too.
\fancypagestyle{plain}{%
  \fancyhf{}%
  \fancyfoot[C]{\stampfootcontent}%
  \renewcommand{\headrulewidth}{0pt}%
  \renewcommand{\footrulewidth}{0.4pt}}
\renewcommand{\stampsource}{\textasciitilde{}/\allowbreak{}prj/\allowbreak{}res/\allowbreak{}36-pmsimstats-ng/\allowbreak{}pmsimstats-ng/\allowbreak{}analysis/\allowbreak{}report/\allowbreak{}15-prognostic-biomarker/\allowbreak{}report.Rmd}
\renewcommand{\stamptime}{2026-11-11 16:14 PST}
\renewcommand{\stampversion}{\detokenize{whitepapers-pre-series-177-gb995bc649-wip}}
\usepackage{amsmath}
\usepackage{amssymb}
\usepackage{booktabs}
\usepackage{longtable}
\usepackage{float}
\usepackage[mathlines]{lineno}
\linenumbers
\usepackage{booktabs}
\usepackage{longtable}
\usepackage{array}
\usepackage{multirow}
\usepackage{wrapfig}
\usepackage{float}
\usepackage{colortbl}
\usepackage{pdflscape}
\usepackage{tabu}
\usepackage{threeparttable}
\usepackage{threeparttablex}
\usepackage[normalem]{ulem}
\usepackage{makecell}
\usepackage{xcolor}
\usepackage{bookmark}
\IfFileExists{xurl.sty}{\usepackage{xurl}}{} % add URL line breaks if available
\urlstyle{same}
% fallback for those not using the hyperref driver hyperxmp:
\makeatletter
\@ifundefined{xmpquote}{\newcommand{\xmpquote}[1]{#1}}{}
\makeatother
\hypersetup{
  pdftitle={When the biomarker is also prognostic: bias in the biomarker-treatment interaction test of aggregated N-of-1 trials and two remedies},
  pdfauthor={pmsimstats team},
  hidelinks,
  pdfcreator={LaTeX via pandoc}}

\title{When the biomarker is also prognostic: bias in the biomarker-treatment interaction test of aggregated N-of-1 trials and two remedies}
\usepackage{etoolbox}
\makeatletter
\providecommand{\subtitle}[1]{% add subtitle to \maketitle
  \apptocmd{\@title}{\par {\large #1 \par}}{}{}
}
\makeatother
\subtitle{First draft, assembled from the material that left paper 02}
\author{pmsimstats team}
\date{November 11, 2026}

\begin{document}
\maketitle

{
\setcounter{tocdepth}{2}
\tableofcontents
}
\listoftables
\setstretch{1.1}
\section*{Abstract}\label{abstract}
\addcontentsline{toc}{section}{Abstract}

\textbf{Background.} The biomarker-treatment interaction test of an
aggregated N-of-1 trial is built for a predictive biomarker. A
biomarker can also be prognostic. It can then be associated with the
natural history, the placebo response or the baseline level. We asked
when such an association biases the interaction test, and what
removes the bias.

\textbf{Methods.} We simulated the Hybrid design under the extended
data-generating process (DGP) of the companion DGP paper (paper 01), with the
parameter values published by Hendrickson et al.~Five prognostic
pathways were examined. Three are correlations of the biomarker with a
component. Two are mean shifts of the natural-history component, one
constant and one growing over the trial. We compared analyses with
and without the design-expectancy terms \(\eta_t\) and \(B^c_i\eta_t\).
A second stress test, run with the driver of paper 01,
fitted ten analyses to trials with four prognostic pathways at
strength 0.3. Its analyses keep the baseline visit as an outcome row.

\textbf{Results.} Under a correlation with the placebo-response component
of 0.45, the analyses without expectancy terms rejected a true null at
0.092 (G9) at \(N = 280\). With the terms, the rate was
0.060 (F1). The two-MCSE band at 1,000 replicates
was \([0.036, 0.064]\).
Under a natural-history shift that grows over the trial, G9 rejected
at 0.073 and F1 at 0.052 at \(N = 70\), with no
baseline row in the analysis. Without a prognostic effect, F1A0
reached 98.2 percent of the power of G10 in the
Hybrid cells. In the second stress test, the analyses with expectancy
terms stayed inside the two-MCSE band under every pathway. Every
analysis without them, G10 included, exceeded the band under at least
one pathway. The bias from the growing shift changed sign when a
baseline row was added. The power loss once attributed to a strong
prognostic correlation came from repaired covariance matrices. Under
the current stop on infeasible matrices, all 10 of
those cells were refused.

\textbf{Conclusions.} For a biomarker that may be prognostic, add
\(\eta_t\) and \(B^c_i\eta_t\) to the recommended analysis of the
companion carryover paper (paper 02). A biomarker that may predict the rate of
natural-history change needs a biomarker-by-time term. That remedy
has not yet been simulated.

\section*{Notation and glossary}\label{notation-and-glossary}
\addcontentsline{toc}{section}{Notation and glossary}

Symbols follow \texttt{analysis/report/NOTATION.md}. Two symbols are not yet
listed there and are marked \emph{proposed}.

{\def\LTcaptype{none} % do not increment counter
\begin{longtable}[]{@{}
  >{\raggedright\arraybackslash}p{(\linewidth - 2\tabcolsep) * \real{0.5000}}
  >{\raggedright\arraybackslash}p{(\linewidth - 2\tabcolsep) * \real{0.5000}}@{}}
\toprule\noalign{}
\begin{minipage}[b]{\linewidth}\raggedright
Symbol
\end{minipage} & \begin{minipage}[b]{\linewidth}\raggedright
Meaning
\end{minipage} \\
\midrule\noalign{}
\endhead
\bottomrule\noalign{}
\endlastfoot
\(B_i\), \(B^c_i\) & biomarker of participant \(i\); centered at the trial's baseline mean \\
\(z_i\) & standardized biomarker \\
\(w_t\) & calendar week of visit \(t\) \\
\(c_{bm}\) & on-drug biomarker-response correlation, the interaction under the correlation encoding \\
\(c_{bm,PB}\), \(c_{bm,TV}\), \(c_{bm,BL}\) & prognostic correlations of the biomarker with the placebo-response component, the natural-history component and the baseline level \\
\(b_{TV}\) (\emph{proposed}) & constant natural-history mean shift per SD of the biomarker \\
\(b_{TV}^{\mathrm{tr}}\) (\emph{proposed}) & coefficient of the trend pathway: a shift of \(b_{TV}^{\mathrm{tr}}\,\sigma_{TV}\,z_i\,g(w_t)\), with \(g\) the natural-history mean course scaled to peak 1 \\
\(\sigma_{BR}\), \(\sigma_{PB}\), \(\sigma_{TV}\) & SDs of the drug-response, placebo-response and natural-history components \\
\(t_{1/2}\), \(t^a_{1/2}\) & carryover half-life of the DGP; half-life assumed by the analysis \\
\(D_{bc,it}\) & exposure-weighted drug indicator built at \(t^a_{1/2}\) \\
\(\eta_t\) & design expectancy: 1 open label, 0.5 blinded, 0 at baseline \\
\(\beta_{bm:D}\) & interaction coefficient, the target of the test \\
\(\beta_\eta\), \(\beta_{bm:\eta}\) & coefficients of \(\eta_t\) and \(B^c_i\eta_t\) \\
\(\kappa\) & mean model standard error divided by the empirical SD of the estimates \\
\(\hat\pi\), MCSE & estimated rejection rate and its Monte Carlo standard error \\
\end{longtable}
}

\begin{itemize}
\tightlist
\item
  \textbf{Correlation encoding.} The interaction is a correlation between
  the biomarker and the drug response. It is present on drug and
  decays after discontinuation at the carryover half-life. Code
  labels it \texttt{mvn}.
\item
  \textbf{Original DGP; extended DGP.} The original DGP is the code
  published by Hendrickson et al.~at commit \texttt{58b32a9}
  (\citeproc{ref-hendrickson2020}{1}). The extended DGP is the one recommended in
  paper 01 (\citeproc{ref-pmsimstats-paper01}{2}), which calls it the recommended DGP.
  It has a decaying biomarker-response
  correlation, the same first-order autoregressive correlation for
  every component with a constant between-component correlation at
  every lag (\(\rho = 0.7\), \(c_1 = 0.2\)), carryover that decays from
  the last on-drug mean, and a stop on correlation matrices that are
  not positive definite. Code and data label it \texttt{refined}.
\item
  \textbf{Published parameter set; simcore2026 set.} The published set is
  the component means and SDs of Hendrickson et al.~Paper 01 calls
  them the default values. The simcore2026
  set is an earlier set used by the paper 02 drivers before
  2026-11-03. Runs under it are labeled historical in this draft.
\item
  \textbf{Predictive biomarker.} One that changes the drug response
  (\(c_{bm} \ne 0\)). \textbf{Prognostic biomarker:} one associated with the
  outcome through another channel. A biomarker can be both.
\item
  \textbf{Prognostic pathway.} One channel of prognostic association. The
  stored data labels are \texttt{prog-BL} (\(c_{bm,BL}\)), \texttt{prog-TV-const}
  (\(c_{bm,TV}\)), \texttt{prog-PB} (\(c_{bm,PB}\)), \texttt{prog-TV-grow} (\(b_{TV}\),
  constant despite its label) and the block T trend pathway
  (\(b_{TV}^{\mathrm{tr}}\)). The second stress test's growing
  natural-history pathway is the trend pathway under another name.
\item
  \textbf{Baseline row.} An outcome row at week 0 holding the baseline
  level, with \(D_{bc,it} = 0\) and \(\eta_t = 0\). The analyses of the
  second stress test include it. Those of the paper 02 driver do not:
  their first row is the week-4 visit.
\item
  \textbf{Feasibility bound.} The largest value of one correlation, given
  the others, for which the joint correlation matrix is positive
  definite on every path of the design.
\item
  \textbf{Silent repair; repaired cell.} Before 2026-10-26 the drivers
  replaced a matrix that was not positive definite with a nearby one
  that was, without warning. A repaired cell is a cell simulated
  under such a matrix.
\item
  \textbf{G-family.} Analyses with no expectancy terms. \textbf{G8:}
  \(D_{bc,it}\) at the true half-life with CR2, an oracle benchmark and
  not a candidate analysis. \textbf{G9:} \(t^a_{1/2}\) chosen by AIC on the
  model with the interaction, with CR2. \textbf{G10:} \(t^a_{1/2}\) chosen by
  AIC on the model without the interaction, the interaction model
  refitted, with CR2. G10 is the recommended analysis of paper 02
  (\citeproc{ref-pmsimstats-paper02}{3}). G1 and G6 are the binary-indicator analyses,
  model-based and with CR2. Stored codes: \texttt{E2cr2}, \texttt{E7cr2},
  \texttt{E7null\_cr2}, \texttt{E1}, \texttt{E1cr2}.
\item
  \textbf{F-family.} The G3 formula plus \(\eta_t\) and \(B^c_i\eta_t\), with
  CR2. \textbf{F1:} \(t^a_{1/2}\) fixed at the half-life policy value of 0.2
  week. \textbf{F1A:} chosen by AIC with the interaction in the model.
  \textbf{F1A0:} chosen by AIC without it. \textbf{F1A0B:} F1A0 with a
  single-resample parametric-bootstrap calibration that repeats the
  selection.
\item
  \textbf{CR2.} The Bell-McCaffrey bias-reduced cluster-robust standard
  error with Satterthwaite degrees of freedom
  (\citeproc{ref-bell2002bias}{4},\citeproc{ref-pustejovsky2018}{5}).
\item
  \textbf{Two-MCSE band.} The interval \(0.05 \pm 2\sqrt{0.05 \times
  0.95/n}\) at \(n\) replicates.
\item
  \textbf{Designs.} Hybrid: open-label titration, then blinded
  discontinuation, then a brief crossover. OL+BDC: open-label
  titration then blinded discontinuation. CO: a two-period crossover.
\item
  \textbf{Single-resample bootstrap.} One bootstrap resample per Monte
  Carlo replicate, with the critical value taken from the pooled
  resamples of the cell.
\end{itemize}

\section{Introduction}\label{sec-intro}

An aggregated N-of-1 trial gives each participant repeated periods on
and off drug. Pooling the participants allows a test of whether a
baseline biomarker modifies the drug effect. That test targets a
predictive biomarker. The interaction test in a parallel-group trial
already needs a large sample (\citeproc{ref-brookes2004}{6}). The N-of-1 design
reduces that cost by comparing each participant with themselves
(\citeproc{ref-senn2002}{7},\citeproc{ref-araujo2016}{8}).

A biomarker can also be prognostic. It may track the severity at
baseline, the natural course of the illness, or the response to
placebo. In a design whose expectancy varies between open-label and
blinded phases, the on-drug and off-drug visits differ in more than
drug exposure. A prognostic association that differs between those
visits can then project onto the interaction term. The test would
then report moderation of the drug effect where there is none. A
prognostic effect is known to bias tests of a
biomarker-by-treatment interaction when the analysis model omits or
misrepresents it (\citeproc{ref-haller2018}{9},\citeproc{ref-kent2020path}{10}).

The companion paper on carryover (\citeproc{ref-pmsimstats-paper02}{3}) studies a
predictive-only biomarker. It recommends G10. That analysis uses an
exposure-weighted drug indicator, chooses its half-life by AIC on the
model without the interaction, and tests with CR2. This paper asks
what changes when the biomarker is also prognostic. It examines five
pathways. It compares the G-family with the F-family, which adds the
design-expectancy terms \(\eta_t\) and \(B^c_i\eta_t\). It also reports
a finding about the simulation itself. An apparent loss of power
under a strong prognostic correlation came from the silent repair of
infeasible covariance matrices. It was not a property of any
analysis.

Paper 01 defines the extended DGP and its prognostic pathways
(\citeproc{ref-pmsimstats-paper01}{2}). It assumes that the biomarker is not
prognostic and refers that case to this paper. Two of its results
therefore appear here. The first is the joint feasibility condition.
A prognostic correlation shares one feasibility condition with
\(c_{bm}\), and the extended DGP refuses a matrix that violates it
(Section \ref{sec-feasibility}, Appendix \ref{sec-jointfeas}). The
second is the prognostic cells of paper 01's stress test, its study
2. That test fitted ten analyses, each with a baseline row, under
four prognostic pathways (Section \ref{sec-tenan}). The studies are
exploratory.

\section{Methods}\label{sec-methods}

\subsection{Data-generating process}\label{sec-dgp}

All published-set cells use the extended DGP under the correlation
encoding, on the Hybrid design, with \(t_{1/2} = 0.5\) week unless
stated. The component means and SDs are those of Hendrickson et al.
Table \ref{tab:tab-params} reads them from the stored run metadata.
The drug-response SD is \(\sigma_{BR} =\) 8. The true
interaction at \(c_{bm}\) is therefore \(-c_{bm}\sigma_{BR}\) per SD of
the biomarker: -3.60 at \(c_{bm} = 0.45\) and
-2.40 at \(c_{bm} = 0.30\). The components are reductions in
symptom severity, so an increase in any component lowers the
outcome. Interaction coefficients are negative.

\begin{table}
\centering
\caption{\label{tab:tab-params}Parameter values of the published-set runs, read from
      the stored run metadata (\texttt{meta.rds}, field
      \texttt{params}). For the
      response components the mean column is the maximum $m$ of the
      modified Gompertz curve.}
\centering
\fontsize{9}{11}\selectfont
\begin{tabular}[t]{lllll}
\toprule
Component & Maximum or mean & SD & Rate (per week) & Displacement\\
\midrule
Natural history (TV) & 6.51 & 10 & 0.35 & 5\\
Placebo response (PB) & 6.51 & 10 & 0.35 & 5\\
Drug response (BR) & 10.99 & 8 & 0.42 & 5\\
Biomarker & 124.3 & 15.4 & -- & --\\
Baseline level & 83.1 & 18.5 & -- & --\\
\bottomrule
\end{tabular}
\end{table}

\subsection{Prognostic pathways}\label{sec-pathways}

A biomarker may be associated with the outcome through channels other
than the drug response. Three pathways are entries in the biomarker
row of the joint correlation matrix. Each sets one correlation with
the others at zero. The first is a correlation with the baseline
level, \(c_{bm,BL}\). The second is a correlation with the
natural-history component at every visit, \(c_{bm,TV}\). The third is a
correlation with the placebo-response component at every visit,
\(c_{bm,PB}\). Their strengths are bounded by feasibility (Section
\ref{sec-feasibility}).

Two further pathways are mean shifts. They are not bounded by the
correlation matrix. The constant shift adds \(b_{TV}\,\sigma_{TV}\,z_i\)
to the natural-history mean at every visit. It is stored as
\texttt{prog-TV-grow} although it does not grow. It is absorbed exactly by
the biomarker main effect of every analysis here (Appendix
\ref{sec-absorb}). It serves only as a check of that absorption. The
trend pathway adds \(b_{TV}^{\mathrm{tr}}\,\sigma_{TV}\,z_i\,g(w_t)\),
where \(g\) is the natural-history mean course scaled to peak 1. It
describes a biomarker that predicts the rate of natural-history
change. A model with no biomarker-by-time term cannot absorb it.

\subsection{Analyses}\label{sec-analyses}

All analyses fit a linear mixed model with a participant random
intercept and continuous-time AR(1) residual correlation
(\texttt{corCAR1}). The G-family model is
\[Y_{it} = \beta_0 + \beta_B B^c_i + \beta_T w_t + \beta_D D_{bc,it}
  + \beta_{bm:D} B^c_i D_{bc,it} + u_i + \epsilon_{it}.\]
\(D_{bc,it}\) is one on drug, \(2^{-t_{sd}/t^a_{1/2}}\) off drug after
first exposure, and zero before first exposure. G8 sets
\(t^a_{1/2} = t_{1/2}\). It is an oracle benchmark and not a candidate
analysis. G9 chooses \(t^a_{1/2}\) by maximum-likelihood AIC over
\(\{0.25, 0.5, 1, 2\}\) weeks with the interaction in the model
(\citeproc{ref-akaike1974new}{11}). G10 makes the same choice on the model without the
interaction and refits the interaction model by REML. The selection
of G9 depends on the coefficient under test. That is a
post-selection inference risk (\citeproc{ref-berk2013valid}{12},\citeproc{ref-leeb2005}{13}). All three
test \(\beta_{bm:D}\) with CR2 and Satterthwaite degrees of freedom.

The F-family adds the design-expectancy terms:
\[Y_{it} = \ldots + \beta_\eta \eta_t + \beta_{bm:\eta} B^c_i \eta_t.\]
The regressors are \texttt{De} and \texttt{bmc:De}. The placebo-response amplitude
scales with \(\eta_t\), and \(\eta_t\) differs between on-drug and
off-drug visits in the Hybrid design. A biomarker correlated with the
placebo response therefore acquires an on/off slope difference by
construction. The term \(B^c_i\eta_t\) absorbs it. F1 fixes
\(t^a_{1/2}\) at 0.2 week. F1A selects it by AIC with the interaction.
F1A0 selects it without the interaction, as G10 does. F1A0B
calibrates F1A0 by a single-resample parametric bootstrap under the
null that repeats the selection. Its candidate grid is \(\{0, 0.2,
0.5, 1, 2\}\) weeks. \(\eta_t\) is constant within every path of the
crossover design. The F-family is therefore not estimable under CO.
F1A0 and G10 differ only in the expectancy terms. That is the
property that lets papers 02 and 15 give one recommendation
(Section \ref{sec-recommend}).

\subsection{Performance measures}\label{sec-perf}

We report the rejection rate \(\hat\pi\) at \(\alpha = 0.05\), its MCSE,
the mean interaction estimate, and \(\kappa\) (\citeproc{ref-morris2019simulation}{14}).
Under the null \(\hat\pi\) is the Type I error. A cell's rejection rate
is compared with its two-MCSE band. That band is a check on gross
departures and not a demonstration of nominal size. Within a job,
cells at the same design and \(N\) share random-number streams, and
every analysis is fitted to the same trials. Contrasts between
analyses are paired, and their standard errors come from the
per-replicate rejection indicators. Comparisons across jobs are
unpaired. Results from different simulation hosts are not
draw-for-draw identical.

\subsection{Stored runs and their parameter sets}\label{sec-sources}

Table \ref{tab:tab-sources} lists every stored run this draft reads.
The second stress test of Section \ref{sec-tenan} comes from another
driver, the stress test of paper 01.
The published-set runs are the reruns of \texttt{docs/74}. The historical
runs used the simcore2026 parameter set and silent repair. They are
kept only as the exhibit of Section \ref{sec-artifact}
(\texttt{docs/74}, decision D5). No historical number enters a result about
an analysis.

\begin{table}
\centering
\caption{\label{tab:tab-sources}Stored runs read by this draft. The parameter set is the
    value recorded in each run's metadata. Runs made before the
    parameter switch of 2026-11-03 record none and used the
    simcore2026 set (docs/74). The second stress test writes no
    metadata file; its dates span its output files, from the first run
    to the G10 fits. Its driver holds the published means and SDs as
    constants. Two analytic files of the historical folder,
    \texttt{pd-check.rds} and \texttt{coupling-ceilings.rds}, are
    read as well. The first depends on the simcore2026 SDs. The
    second depends on the correlation matrix only.}
\centering
\resizebox{\ifdim\width>\linewidth\linewidth\else\width\fi}{!}{
\fontsize{8}{10}\selectfont
\begin{tabular}[t]{lllll}
\toprule
Run & Folder under analysis/data/ & Parameter set recorded & Date & Status\\
\midrule
Prognostic stress test & \texttt{carryover-sensitivity/published/g9-prognostic-stress} & published & 2026-11-04 & published set\\
Feasible grid and large N & \texttt{carryover-sensitivity/published/g9-gap-fill/C} & published & 2026-11-05 & published set\\
Trend pathway (block T) & \texttt{carryover-sensitivity/published/referee-variants/block-T} & published & 2026-11-04 & published set\\
Hybrid nulls (block N) & \texttt{carryover-sensitivity/published/referee-variants/block-N} & published & 2026-11-04 & published set\\
Hybrid alternatives (block A) & \texttt{carryover-sensitivity/published/referee-variants/block-A} & published & 2026-11-04 & published set\\
\addlinespace
Original 30-cell grid & \texttt{carryover-sensitivity/published/tier1-g-vs-f-e0} & published & 2026-11-04 & published set\\
Scenario cells & \texttt{carryover-sensitivity/published/g-vs-f-comparison} & published & 2026-11-04 & published set\\
Second stress test, three runs and post hoc fits & \texttt{quick-sim/carryover-closed-form/strawman-stress} & driver constants & 2026-10-07 to 2026-11-07 & published set\\
Joint feasibility bounds & \texttt{quick-sim/carryover-closed-form/rho-stage/prognostic-joint-ceilings.csv} & correlations only & 2026-10-27 & analytic\\
Prognostic stress test & \texttt{carryover-sensitivity/g9-prognostic-stress} & not recorded & 2026-10-26 & historical, repaired\\
\addlinespace
Random-slope variants & \texttt{carryover-sensitivity/g9-random-effects} & not recorded & 2026-10-26 & historical, repaired\\
\bottomrule
\end{tabular}}
\end{table}

\section{Results}\label{sec-results}

\subsection{Feasibility of a prognostic correlation}\label{sec-feasibility}

A prognostic correlation occupies the same biomarker row of the
correlation matrix as \(c_{bm}\). The two therefore share one
feasibility condition. Writing the row as a sum of fixed patterns
scaled by the correlations, the condition becomes a quadratic form in
the correlations that must be below one. The largest value of one
correlation given the others is the positive root of a quadratic
(Appendix \ref{sec-jointfeas}). Table \ref{tab:tab-jointceil} gives
the bounds for the Hybrid design under the extended DGP. They depend
on the correlation matrix only, not on the component SDs, so they
hold for both parameter sets.

\begin{table}
\centering
\caption{\label{tab:tab-jointceil}Largest feasible prognostic correlation given the
      biomarker-response correlation $c_{bm}$, and the reverse, under
      the extended DGP (separable AR(1), $\rho = 0.7$, $c_1 = 0.2$),
      Hybrid design, minimum over its four paths. Correlation: the
      quantity whose bound is given. $c_{bm,PB}$, $c_{bm,TV}$,
      $c_{bm,BL}$: correlation of the biomarker with the
      placebo-response component, the natural-history component and
      the baseline level. None: no value is feasible. Computed from
      the correlation matrix, no simulation
      (\texttt{18-prognostic-joint-ceiling.R}).}
\centering
\fontsize{9}{11}\selectfont
\begin{tabular}[t]{llll}
\toprule
Correlation & Given & $t_{1/2} = 0.2$ week & $t_{1/2} = 0.5$ week\\
\midrule
$c_{bm,PB}$ or $c_{bm,TV}$ & alone & 0.513 & 0.513\\
 & $c_{bm} = 0.25$ & 0.466 & 0.477\\
 & $c_{bm} = 0.30$ & 0.435 & 0.451\\
 & $c_{bm} = 0.45$ & 0.243 & 0.306\\
$c_{bm,BL}$ & alone & 1.000 & 1.000\\
 & $c_{bm} = 0.25$ & 0.857 & 0.876\\
 & $c_{bm} = 0.30$ & 0.786 & 0.815\\
 & $c_{bm} = 0.45$ & 0.374 & 0.496\\
$c_{bm}$ & alone & 0.485 & 0.518\\
 & $c_{bm,PB} = 0.3$ & 0.422 & 0.454\\
 & $c_{bm,PB} = 0.6$ & none & none\\
\bottomrule
\end{tabular}
\end{table}

At this paper's half-life of 0.5 week, a prognostic correlation of
0.3 lies inside the feasible region at every \(c_{bm}\) up to 0.45, by
0.006 at 0.45. At 0.2 week it is feasible only up to
\(c_{bm} =\) 0.42. A prognostic correlation of 0.6 with the
placebo-response or natural-history component is infeasible with any
\(c_{bm}\). That is the combination that the repaired runs of Section
\ref{sec-artifact} simulated. The correlations with the
placebo-response and the natural-history components have identical
bounds, because their patterns are the same constant on every visit
of one component. A correlation with the baseline level has a looser
bound. The baseline level is uncorrelated with the response
components and enters only through its square. Outside the table,
the OL+BDC bounds are within 0.034 of the Hybrid
values at the same half-life. The crossover bounds are looser on
\(c_{bm}\) and \(c_{bm,BL}\), and slightly tighter on \(c_{bm,PB}\) alone
(0.487 against 0.513). Each bound was checked
by the sign of the smallest eigenvalue of the full matrix: positive
at 0.999 times the bound and negative at 1.001 times it. The feasible
grid of Section \ref{sec-grid} was chosen from these bounds.

\subsection{Correlation pathways at strength 0.3}\label{sec-stress}

We crossed the three correlation pathways and the constant shift
with \(c_{bm} \in \{0, 0.45\}\) and \(N \in \{35, 70\}\), at
500 replicates per cell. Table \ref{tab:tab-stress} gives
the rejection rates. The two-MCSE band is \([0.031, 0.069]\).

\begin{table}
\centering
\caption{\label{tab:tab-stress}Rejection rate under each prognostic pathway
      at strength 0.3, published parameter set. Hybrid, $t_{1/2} =
      0.5$ week, 500 replicates per cell, the same trials for every
      analysis. Rows with $c_{bm} = 0$ are Type I error (band $[0.031, 0.069]$).
      G8 is the oracle benchmark.}
\centering
\fontsize{9}{11}\selectfont
\begin{tabular}[t]{llrllllll}
\toprule
Pathway & $c_{bm}$ & $N$ & G8 & G9 & G10 & F1 & F1A0 & F1A0B\\
\midrule
$c_{bm,PB}$ & 0.00 & 35 & 0.040 & 0.044 & 0.038 & 0.032 & 0.034 & 0.032\\
 & 0.00 & 70 & 0.060 & 0.058 & 0.058 & 0.042 & 0.038 & 0.032\\
 & 0.45 & 35 & 0.484 & 0.500 & 0.468 & 0.380 & 0.379 & 0.415\\
 & 0.45 & 70 & 0.786 & 0.792 & 0.780 & 0.658 & 0.678 & 0.646\\
$c_{bm,TV}$ & 0.00 & 35 & 0.038 & 0.042 & 0.038 & 0.034 & 0.032 & 0.034\\
 & 0.00 & 70 & 0.056 & 0.050 & 0.048 & 0.042 & 0.036 & 0.030\\
 & 0.45 & 35 & 0.418 & 0.431 & 0.402 & 0.371 & 0.377 & 0.405\\
 & 0.45 & 70 & 0.714 & 0.716 & 0.706 & 0.662 & 0.668 & 0.654\\
$c_{bm,BL}$ & 0.00 & 35 & 0.040 & 0.044 & 0.038 & 0.032 & 0.032 & 0.034\\
 & 0.00 & 70 & 0.056 & 0.054 & 0.048 & 0.042 & 0.038 & 0.032\\
 & 0.45 & 35 & 0.426 & 0.448 & 0.414 & 0.381 & 0.377 & 0.411\\
 & 0.45 & 70 & 0.720 & 0.730 & 0.712 & 0.682 & 0.692 & 0.648\\
$b_{TV}$ (constant) & 0.00 & 35 & 0.038 & 0.042 & 0.038 & 0.030 & 0.032 & 0.034\\
 & 0.00 & 70 & 0.056 & 0.054 & 0.048 & 0.042 & 0.038 & 0.032\\
 & 0.45 & 35 & 0.409 & 0.424 & 0.393 & 0.363 & 0.369 & 0.403\\
 & 0.45 & 70 & 0.720 & 0.724 & 0.710 & 0.694 & 0.698 & 0.680\\
\bottomrule
\end{tabular}
\end{table}

The G9 Type I error was 0.042 to 0.058 across the six null cells
of the three correlation pathways. G10 gave 0.038 to 0.058 and
F1A0B 0.030 to 0.034. Every G9 and G10 null rate was inside the
band. The G9 power at \(c_{bm} = 0.45\) was 0.424 to 0.792. The
no-confound reference from block A of the referee-variant runs is
0.442 at \(N = 35\) and 0.728 at \(N = 70\). The
largest difference from it was +0.064, under the
placebo-response pathway. Section \ref{sec-grid} shows that this
excess is bias in the estimate. That comparison is
unpaired, against 1,000-replicate cells of another job. F1A0B
differed from G10 in power by -0.134 to +0.010 on the same trials.
The constant shift gave the same rejection at strengths 0.3 and 0.6
in every analysis. In the scenario job, the constant-shift null cell
differs from the unshifted Hybrid null cell by at most
0.0000 in any analysis. That is the absorption of Appendix
\ref{sec-absorb}.

A longer run at the placebo-response pathway sharpens the size
comparison. Block N of the referee-variant runs fitted G8, G9, G10
and F1 to 5,000 replicates in each of 4
Hybrid null cells at \(c_{bm,PB} =\) 0.3
(Table \ref{tab:tab-blockN}). The band is \([0.044, 0.056]\). The
paired difference from G8 was +0.0056 to +0.0134 for G9, -0.0012 to +0.0016
for G10 and -0.0072 to -0.0028 for F1. G9 was above G8 in every cell,
by at least 3.5 paired standard errors. G10 was within
two paired standard errors of G8 in every cell. The number of cells
above the band was 2 for G8, 4 for
G9, 2 for G10 and 0 for F1.

\begin{table}
\centering
\caption{\label{tab:tab-blockN}Block N, published parameter set: null
      rejection rate (MCSE) under the placebo-response pathway at
      $c_{bm,PB} = 0.3$, Hybrid, 5,000 replicates per cell, the same
      trials for every analysis, with the paired difference from G8
      and its paired SE. Band $[0.044, 0.056]$.}
\centering
\fontsize{8}{10}\selectfont
\begin{tabular}[t]{lllll}
\toprule
Cell & Method & Rejection & Diff. vs G8 (SE) & $\kappa$\\
\midrule
$t_{1/2} = 0.5$, $N = 35$ & G8 & 0.0566 (0.0033) & -- & 0.95\\
 & G9 & 0.0653 (0.0035) & +0.0086 (0.0017) & 0.88\\
 & G10 & 0.0573 (0.0033) & +0.0006 (0.0016) & 0.94\\
 & F1 & 0.0541 (0.0032) & -0.0028 (0.0027) & 0.96\\
$t_{1/2} = 0.5$, $N = 70$ & G8 & 0.0530 (0.0032) & -- & 1.00\\
 & G9 & 0.0586 (0.0033) & +0.0056 (0.0016) & 0.96\\
 & G10 & 0.0538 (0.0032) & +0.0008 (0.0015) & 1.00\\
 & F1 & 0.0474 (0.0030) & -0.0056 (0.0025) & 1.00\\
$t_{1/2} = 1.0$, $N = 35$ & G8 & 0.0469 (0.0030) & -- & 0.97\\
 & G9 & 0.0601 (0.0034) & +0.0130 (0.0021) & 0.88\\
 & G10 & 0.0457 (0.0030) & -0.0012 (0.0019) & 0.95\\
 & F1 & 0.0411 (0.0028) & -0.0060 (0.0030) & 0.98\\
$t_{1/2} = 1.0$, $N = 70$ & G8 & 0.0606 (0.0034) & -- & 0.95\\
 & G9 & 0.0740 (0.0037) & +0.0134 (0.0022) & 0.89\\
 & G10 & 0.0622 (0.0034) & +0.0016 (0.0021) & 0.93\\
 & F1 & 0.0534 (0.0032) & -0.0072 (0.0033) & 0.96\\
\bottomrule
\end{tabular}
\end{table}

\subsection{The feasible prognostic grid}\label{sec-grid}

The supplementary study ran a grid of correlations chosen inside the
feasibility bounds. At \(c_{bm} = 0.30\) it set \(c_{bm,PB}\) and
\(c_{bm,TV}\) at 0.40 and 0.45, and \(c_{bm,BL}\) at 0.6 and 0.8. It also
ran \(c_{bm,BL} = 0.45\) at \(c_{bm} = 0.45\). Each cell was run with and
without a true interaction, at \(N = 70\) and 500 replicates.
The strongest \(c_{bm,PB}\) lies 0.001 below its bound at
\(c_{bm} = 0.30\). The strongest \(c_{bm,BL}\) lies 0.015
below its bound. Table \ref{tab:tab-grid} gives the rejection rates
and Table \ref{tab:tab-grid-est} the mean estimates.

\begin{table}
\centering
\caption{\label{tab:tab-grid}Rejection rate on the feasible prognostic grid,
      published parameter set. Hybrid, $N = 70$, $t_{1/2} = 0.5$
      week, 500 replicates per cell. Rows with $c_{bm} = 0$ are Type I
      error (band $[0.031, 0.069]$). The others are power.}
\centering
\fontsize{9}{11}\selectfont
\begin{tabular}[t]{lllllll}
\toprule
Cell & G8 & G9 & G10 & F1 & F1A0 & F1A0B\\
\midrule
$c_{bm,PB}$ = 0.40, $c_{bm} = 0.00$ & 0.066 & 0.056 & 0.054 & 0.060 & 0.058 & 0.060\\
$c_{bm,PB}$ = 0.40, $c_{bm} = 0.30$ & 0.564 & 0.564 & 0.540 & 0.400 & 0.382 & 0.384\\
$c_{bm,PB}$ = 0.45, $c_{bm} = 0.00$ & 0.066 & 0.058 & 0.054 & 0.058 & 0.058 & 0.058\\
$c_{bm,PB}$ = 0.45, $c_{bm} = 0.30$ & 0.582 & 0.580 & 0.566 & 0.380 & 0.390 & 0.370\\
$c_{bm,TV}$ = 0.40, $c_{bm} = 0.00$ & 0.040 & 0.052 & 0.048 & 0.054 & 0.052 & 0.056\\
$c_{bm,TV}$ = 0.40, $c_{bm} = 0.30$ & 0.418 & 0.432 & 0.412 & 0.408 & 0.384 & 0.386\\
$c_{bm,TV}$ = 0.45, $c_{bm} = 0.00$ & 0.040 & 0.054 & 0.050 & 0.054 & 0.050 & 0.050\\
$c_{bm,TV}$ = 0.45, $c_{bm} = 0.30$ & 0.439 & 0.438 & 0.410 & 0.406 & 0.388 & 0.392\\
$c_{bm,BL}$ = 0.45, $c_{bm} = 0.00$ & 0.046 & 0.066 & 0.062 & 0.060 & 0.056 & 0.056\\
$c_{bm,BL}$ = 0.45, $c_{bm} = 0.45$ & 0.742 & 0.746 & 0.740 & 0.732 & 0.726 & 0.724\\
$c_{bm,BL}$ = 0.60, $c_{bm} = 0.00$ & 0.046 & 0.066 & 0.062 & 0.060 & 0.054 & 0.060\\
$c_{bm,BL}$ = 0.60, $c_{bm} = 0.30$ & 0.432 & 0.453 & 0.427 & 0.399 & 0.383 & 0.391\\
$c_{bm,BL}$ = 0.80, $c_{bm} = 0.00$ & 0.045 & 0.061 & 0.057 & 0.057 & 0.055 & 0.061\\
$c_{bm,BL}$ = 0.80, $c_{bm} = 0.30$ & 0.444 & 0.429 & 0.404 & 0.404 & 0.394 & 0.414\\
\bottomrule
\end{tabular}
\end{table}

\begin{table}
\centering
\caption{\label{tab:tab-grid-est}Mean interaction estimate on the feasible prognostic
      grid (same cells as the rejection table). The true value is
      $-c_{bm}\sigma_{BR}$ per SD of the biomarker.}
\centering
\fontsize{9}{11}\selectfont
\begin{tabular}[t]{llllll}
\toprule
Cell & Truth & G9 & G10 & F1 & F1A0\\
\midrule
$c_{bm,PB}$ = 0.40, $c_{bm} = 0.00$ & -0.00 & -0.49 & -0.47 & -0.09 & -0.06\\
$c_{bm,PB}$ = 0.40, $c_{bm} = 0.30$ & -2.40 & -2.87 & -2.78 & -2.17 & -2.32\\
$c_{bm,PB}$ = 0.45, $c_{bm} = 0.00$ & -0.00 & -0.54 & -0.52 & -0.09 & -0.06\\
$c_{bm,PB}$ = 0.45, $c_{bm} = 0.30$ & -2.40 & -2.94 & -2.84 & -2.18 & -2.34\\
$c_{bm,TV}$ = 0.40, $c_{bm} = 0.00$ & -0.00 & -0.08 & -0.09 & -0.09 & -0.05\\
$c_{bm,TV}$ = 0.40, $c_{bm} = 0.30$ & -2.40 & -2.45 & -2.39 & -2.17 & -2.32\\
$c_{bm,TV}$ = 0.45, $c_{bm} = 0.00$ & -0.00 & -0.08 & -0.08 & -0.08 & -0.04\\
$c_{bm,TV}$ = 0.45, $c_{bm} = 0.30$ & -2.40 & -2.47 & -2.41 & -2.19 & -2.35\\
$c_{bm,BL}$ = 0.45, $c_{bm} = 0.00$ & -0.00 & -0.10 & -0.09 & -0.09 & -0.06\\
$c_{bm,BL}$ = 0.45, $c_{bm} = 0.45$ & -3.60 & -3.64 & -3.56 & -3.22 & -3.47\\
$c_{bm,BL}$ = 0.60, $c_{bm} = 0.00$ & -0.00 & -0.09 & -0.09 & -0.09 & -0.07\\
$c_{bm,BL}$ = 0.60, $c_{bm} = 0.30$ & -2.40 & -2.46 & -2.39 & -2.17 & -2.31\\
$c_{bm,BL}$ = 0.80, $c_{bm} = 0.00$ & -0.00 & -0.09 & -0.09 & -0.11 & -0.08\\
$c_{bm,BL}$ = 0.80, $c_{bm} = 0.30$ & -2.40 & -2.47 & -2.41 & -2.18 & -2.34\\
\bottomrule
\end{tabular}
\end{table}

The G9 Type I error was 0.052 to 0.066 across the
7 null cells, inside the band in every one. G10 gave
0.048 to 0.062 and F1 0.054 to 0.060. The rejection rates
hide a bias that the mean estimates show. Under the placebo-response
pathway with no true interaction, the mean G9 estimate was
-0.54 to -0.49. G10 gave -0.52 to -0.47. F1 gave
-0.09. Under the other two pathways the mean G9
estimate was -0.10 to -0.08. At \(c_{bm} =\) 0.30
the same bias appears as excess power. G9 reached 0.564 to 0.580
under the placebo-response pathway and 0.429 to 0.453 under the
others. Its mean estimate was -2.94 to -2.87 and
-2.47 to -2.45 against a true value of -2.40. F1
reached 0.380 to 0.400 under the placebo-response pathway. Its
mean estimate was -2.19 to -2.17 across the grid. The F1 estimate
is attenuated because its fixed half-life of 0.2 week is below the
true 0.5 week (Section \ref{sec-gvf}). The excess power of the
G-family under the placebo-response pathway is therefore not evidence
of better power. It is the same bias as under the null.

\subsection{Larger samples under the placebo-response pathway}\label{sec-largen}

A bias of a fixed size becomes Type I error inflation as \(N\) grows,
because the standard error shrinks. The supplementary study ran the
placebo-response pathway with no true interaction at \(N \in \{140,
280\}\) and \(c_{bm,PB} \in \{0.30, 0.45\}\), at 1,000
replicates. The band is \([0.036, 0.064]\). Table \ref{tab:tab-largen}
gives the result.

\begin{table}
\centering
\caption{\label{tab:tab-largen}Type I error under the placebo-response
      pathway at larger samples, published parameter set. Hybrid,
      $t_{1/2} = 0.5$ week, $c_{bm} = 0$, 1,000 replicates at $N = 140$
      and $280$ (band $[0.036, 0.064]$). The $N = 70$ row is from the feasible grid
      at 500 replicates.}
\centering
\fontsize{9}{11}\selectfont
\begin{tabular}[t]{llllllll}
\toprule
Cell & G1 & G8 & G9 & G10 & F1 & F1A0 & F1A0B\\
\midrule
$c_{bm,PB} = 0.30$, $N = 140$ & 0.048 & 0.060 & 0.062 & 0.058 & 0.052 & 0.050 & 0.041\\
$c_{bm,PB} = 0.30$, $N = 280$ & 0.055 & 0.065 & 0.067 & 0.066 & 0.056 & 0.053 & 0.050\\
$c_{bm,PB} = 0.45$, $N = 70$ & 0.046 & 0.066 & 0.058 & 0.054 & 0.058 & 0.058 & 0.058\\
$c_{bm,PB} = 0.45$, $N = 140$ & 0.057 & 0.076 & 0.077 & 0.073 & 0.051 & 0.049 & 0.042\\
$c_{bm,PB} = 0.45$, $N = 280$ & 0.069 & 0.094 & 0.092 & 0.089 & 0.060 & 0.058 & 0.053\\
\bottomrule
\end{tabular}
\end{table}

At \(N = 280\) and \(c_{bm,PB} = 0.45\), G9 rejected at
0.092, G10 at 0.089 and G8 at 0.094.
The model-based binary-indicator analysis G1 rejected at
0.069. The F-family stayed inside the band: F1 at
0.060, F1A0 at 0.058 and F1A0B at
0.053. On the same trials, the paired difference in
rejection of F1 minus G9 was -0.032 (paired SE
0.009). That of F1A0 minus G10 was -0.031
(0.009). The G9 rate rose with \(N\):
0.058 at 70, 0.077 at 140 and
0.092 at 280. Its mean estimate at \(N = 280\) was
-0.36 against +0.05 for F1. The standard
errors were calibrated for both, with \(\kappa\) of 0.96 and
0.97. The excess is therefore bias, not an
underestimated standard error. At \(c_{bm,PB} = 0.30\) the rates were
0.067 for G9 and 0.056 for F1.

\subsection{The trend pathway}\label{sec-trend}

Block T of the referee-variant runs simulated a natural-history shift
that grows over the trial, with \(b_{TV}^{\mathrm{tr}} =\)
0.3. It is a mean shift, so the feasibility bound does
not limit it. G8, G9 and F1 were fitted at \(N = 70\) and
1,000 replicates (Table \ref{tab:tab-trend}). The band is
\([0.036, 0.064]\).

\begin{table}
\centering
\caption{\label{tab:tab-trend}Block T, published parameter set: rejection
      rate (MCSE), mean interaction estimate, empirical SE and
      calibration ratio under the trend pathway
      ($b_{TV}^{\mathrm{tr}} = 0.3$). Hybrid, $N = 70$, $t_{1/2} =
      0.5$ week, 1,000 replicates. The true value is 0 at the null and
      -3.60 at $c_{bm} = 0.45$. Band at the null $[0.036, 0.064]$.}
\centering
\fontsize{9}{11}\selectfont
\begin{tabular}[t]{llllll}
\toprule
Cell & Method & Rejection & Mean estimate & Empirical SE & $\kappa$\\
\midrule
null ($c_{bm} = 0$) & G8 & 0.063 (0.008) & +0.43 & 1.45 & 0.97\\
 & G9 & 0.073 (0.008) & +0.42 & 1.46 & 0.93\\
 & F1 & 0.052 (0.007) & +0.24 & 1.27 & 0.98\\
$c_{bm} = 0.45$ & G8 & 0.629 (0.015) & -3.20 & 1.37 & 0.97\\
 & G9 & 0.631 (0.015) & -3.12 & 1.41 & 0.90\\
 & F1 & 0.632 (0.015) & -2.88 & 1.23 & 0.97\\
\bottomrule
\end{tabular}
\end{table}

The trend pathway biased the G-family. Under the null G8 rejected at
0.063 and G9 at 0.073. Their mean estimates were
+0.43 and +0.42 where the truth is zero. The
G8 empirical SE was 1.45. F1 rejected at 0.052,
with a mean estimate of +0.24. Paired against G8, F1
differed by -0.011 (SE 0.006) and G9 by
+0.010 (0.005). The expectancy terms therefore
absorb part of the shift but not all of it. At \(c_{bm} = 0.45\) the
bias is toward zero. G9 power was 0.631 with a mean
estimate of -3.12. In the no-trend cell of block A it
was 0.728 with -3.47. That comparison is
unpaired. The mechanism is the design. In the Hybrid design
\(D_{bc,it}\) falls with calendar time, so a biomarker that predicts
the trend projects onto the interaction term of a model with no
biomarker-by-time term. The paper 02 driver has no baseline row. The
second stress test applied the same shift to analyses with a
baseline row and found a bias of the opposite sign (Section
\ref{sec-reconcile}).

\subsection{The biomarker-by-time remedy: RESULTS PENDING}\label{sec-bmt}

\fbox{\parbox{0.95\linewidth}{\textbf{RESULTS PENDING.} The
simulation for this section has not been run. This section contains
no results. It will be written from the block described in the split
plan, Section 3.1, once that block is stored.}}

The remedy adds a biomarker-by-time term \(B^c_i w_t\) to the analysis
model, as \texttt{bmc:t}. Its coefficient would be written
\(\beta_{bm:T}\) (\emph{proposed}). A quadratic variant adds \(B^c_i w_t^2\).
The planned block crosses the trend pathway, present or absent, with
a null and an alternative interaction. It fits G8, G9, G10 and F1A0
with and without the term to the same trials. It answers three
questions. Does the term remove the trend bias in size and in the
mean estimate? What does it cost in power when there is no trend?
Does a linear term suffice when the natural-history course is not a
straight line? The driver changes it needs are listed in the
Reproducibility section.

\subsection{Expectancy terms without a prognostic effect}\label{sec-gvf}

The F-family's cost when the biomarker is predictive only was
measured on the original 30-cell grid of paper 02 under the published
parameter set (500 replicates per cell). The F-family is
not estimable under CO, so only the Hybrid and OL+BDC cells enter.
Table \ref{tab:tab-gvf} gives the rejection rates.

\begin{table}
\centering
\caption{\label{tab:tab-gvf}Original 30-cell grid of paper 02, published
      parameter set, Hybrid and OL+BDC cells: rejection rate with no
      prognostic effect. 500 replicates per cell, the same trials for
      every analysis. Rows with $c_{bm} = 0$ are Type I error (band
      $[0.031, 0.069]$).}
\centering
\fontsize{9}{11}\selectfont
\begin{tabular}[t]{lllrlllll}
\toprule
Design & $c_{bm}$ & $t_{1/2}$ & $N$ & G9 & G10 & F1 & F1A0 & F1A0B\\
\midrule
Hybrid & 0.00 & 1.0 & 35 & 0.078 & 0.060 & 0.054 & 0.056 & 0.062\\
 & 0.00 & 1.0 & 70 & 0.066 & 0.062 & 0.052 & 0.054 & 0.054\\
 & 0.30 & 1.0 & 35 & 0.179 & 0.155 & 0.136 & 0.140 & 0.154\\
 & 0.30 & 1.0 & 70 & 0.327 & 0.299 & 0.280 & 0.280 & 0.272\\
 & 0.45 & 0.0 & 35 & 0.512 & 0.508 & 0.486 & 0.486 & 0.486\\
 & 0.45 & 0.0 & 70 & 0.868 & 0.868 & 0.872 & 0.874 & 0.874\\
 & 0.45 & 0.5 & 35 & 0.394 & 0.376 & 0.362 & 0.374 & 0.360\\
 & 0.45 & 0.5 & 70 & 0.749 & 0.737 & 0.718 & 0.711 & 0.699\\
 & 0.45 & 1.0 & 35 & 0.317 & 0.277 & 0.248 & 0.279 & 0.287\\
 & 0.45 & 1.0 & 70 & 0.597 & 0.553 & 0.501 & 0.561 & 0.543\\
OL+BDC & 0.00 & 1.0 & 35 & 0.052 & 0.044 & 0.064 & 0.056 & 0.048\\
 & 0.00 & 1.0 & 70 & 0.056 & 0.050 & 0.060 & 0.056 & 0.040\\
 & 0.30 & 1.0 & 35 & 0.114 & 0.108 & 0.086 & 0.106 & 0.094\\
 & 0.30 & 1.0 & 70 & 0.204 & 0.198 & 0.150 & 0.192 & 0.140\\
 & 0.45 & 0.0 & 35 & 0.401 & 0.390 & 0.382 & 0.372 & 0.384\\
 & 0.45 & 0.0 & 70 & 0.721 & 0.715 & 0.695 & 0.682 & 0.644\\
 & 0.45 & 0.5 & 35 & 0.305 & 0.289 & 0.220 & 0.258 & 0.270\\
 & 0.45 & 0.5 & 70 & 0.567 & 0.563 & 0.510 & 0.508 & 0.454\\
 & 0.45 & 1.0 & 35 & 0.202 & 0.198 & 0.142 & 0.178 & 0.158\\
 & 0.45 & 1.0 & 70 & 0.385 & 0.365 & 0.300 & 0.340 & 0.296\\
\bottomrule
\end{tabular}
\end{table}

The null rates were 0.044 to 0.062 for G10, 0.054 to 0.056
for F1A0, 0.040 to 0.062 for F1A0B and 0.052 to 0.064 for F1.
Over the alternative cells of the Hybrid design, F1A0 reached
98.2 percent of the mean power of G10, F1A0B
97.4 percent and F1 95.5
percent. Under OL+BDC the figures were 93.3,
86.3 and 88.0 percent. The mean
paired difference from G10 over the 8 Hybrid alternative
cells was -0.009 for F1A0 (SE
0.004) and -0.022 for F1
(0.005). Under OL+BDC it was
-0.024 (0.005) and
-0.043 (0.005). The pooled
standard errors treat the cells as independent. Cells at the same
\(N\) share random-number streams, so they are approximate.

F1 also loses power where its fixed half-life is wrong. In block A
at a true half-life of 2 weeks and \(c_{bm} = 0.45\), F1 reached
0.163 to 0.279 against 0.192 to 0.359 for G10. Its mean estimate
was -1.85 to -1.72 against -3.09 to -2.93 for G10 and a true
value of -3.60. The fixed half-life is the cost, not the
expectancy terms. F1A0, which selects the half-life, is the fair
comparator for the terms themselves.

\begin{table}
\centering
\caption{\label{tab:tab-scen}Scenario cells of paper 02, published parameter
      set, no prognostic effect: rejection rate. Hybrid unless
      stated, $N = 70$, $t_{1/2} = 0.5$ week unless stated, 500
      replicates. In the $t_{1/2} = 2$ week rows the fixed half-life
      of F1 is a tenth of the true value.}
\centering
\fontsize{9}{11}\selectfont
\begin{tabular}[t]{llllll}
\toprule
Cell & G9 & G10 & F1 & F1A0 & F1A0B\\
\midrule
Hybrid, null & 0.060 & 0.052 & 0.068 & 0.058 & 0.056\\
OL+BDC, null & 0.060 & 0.058 & 0.048 & 0.048 & 0.046\\
$t_{1/2} = 2$ wk, null & 0.078 & 0.066 & 0.068 & 0.058 & 0.052\\
Weibull $k = 2$ data, null & 0.062 & 0.058 & 0.068 & 0.064 & 0.054\\
Hybrid, $c_{bm} = 0.45$ & 0.754 & 0.736 & 0.728 & 0.722 & 0.710\\
Hybrid, $c_{bm} = 0.25$ & 0.296 & 0.288 & 0.290 & 0.282 & 0.280\\
OL+BDC, $c_{bm} = 0.45$ & 0.576 & 0.566 & 0.524 & 0.532 & 0.480\\
$t_{1/2} = 2$ wk, $c_{bm} = 0.25$ & 0.162 & 0.138 & 0.104 & 0.142 & 0.132\\
Weibull $k = 2$ data, $c_{bm} = 0.25$ & 0.378 & 0.374 & 0.362 & 0.362 & 0.332\\
\bottomrule
\end{tabular}
\end{table}

\subsection{A second stress test: ten analyses under the extended DGP}\label{sec-tenan}

This section reports the prognostic cells of a second stress test, run
with the driver of paper 01
(\texttt{analysis/scripts/quick-sim/carryover-closed-form/11-strawman-stress-test.R}).
Paper 01 reports the cells of the same study, its study 2, that have
no prognostic pathway (\citeproc{ref-pmsimstats-paper01}{2}). The trials come from the
extended DGP with the published parameter set, in the Hybrid design:
\(N = 70\) with fixed allocation (18, 18, 17, 17), \(t_{1/2} = 0.5\) week,
\(\rho = 0.7\) and \(c_1 = 0.2\). Four null cells (\(c_{bm} = 0\)) each add
one prognostic pathway of strength 0.3: \(c_{bm,TV} = 0.3\),
\(c_{bm,PB} = 0.3\), \(c_{bm,BL} = 0.3\), and a mean shift of the
natural-history component of 0.3 of its SD per SD of the biomarker at
the end of its curve, growing along the curve. The last is labeled the
growing natural-history pathway here. It is not the stored
\texttt{prog-TV-grow} cell of the paper 02 driver, which is constant (Section
\ref{sec-pathways}). It is the trend pathway of Section
\ref{sec-trend}, with the same coefficient (Section
\ref{sec-reconcile}). Each cell has 1,000 replicates, so the
single-cell two-MCSE band is 0.036 to 0.064.

Fourteen analyses were fitted to the same trials, and ten are
reported. They are named by their model in Table
\ref{tab:tab-tenan}, not by the G and F codes of this paper. Every
one of them keeps the baseline visit as an outcome row. The
expectancy-term model adds \(\eta_t\) and \(B^c_i\eta_t\), as the F-family
does. The free-baseline model adds \texttt{bmc:bl0}, a separate biomarker
slope for the baseline row. Analyses with a fixed analysis half-life
use 0.5 week, the true value.

The study was run three times with the same seed rule,
\(20261007 + 1000 \times \text{cell} + \text{chunk}\). The first run
used the compounding carryover rule on the mean and fixed allocation.
The second used carryover decaying from the last on-drug mean and
randomized allocation. The third, reported here, is the extended DGP
as defined: decaying carryover and fixed allocation. It differs from
the first only in the carryover rule and from the second only in the
allocation. Failed fits were excluded from the denominator. They
were at most 0.6\% in any cell. The study is exploratory. Its cells and
analyses followed a working protocol that was revised as results came
in. Two cells with a prognostic association with baseline severity,
and three analyses, were added after the first results.

\begin{table}
\centering
\caption{\label{tab:tab-tenan}Second stress test: null rejection rates under a
    prognostic biomarker ($c_{bm} = 0$). Each pathway has strength
    0.3: a correlation of 0.3 with the natural-history component
    (constant), the placebo-response component or baseline severity,
    and for the growing natural-history pathway a mean shift of the
    natural-history component of 0.3 of its SD (3 points) per SD of
    biomarker at the end of its curve, growing along the curve. 1,000
    replicates per cell (two-MCSE band 0.036 to 0.064 for a single
    cell). Third run: the extended DGP, with carryover decaying from
    the last on-drug mean and fixed path allocation. Every analysis
    has a baseline row. Rates above 0.075 in bold. The first run is
    in Appendix \ref{sec-firstrun}.}
\centering
\fontsize{9}{11}\selectfont
\begin{tabular}[t]{>{\raggedright\arraybackslash}p{5.5cm}>{\raggedright\arraybackslash}p{2.1cm}>{\raggedright\arraybackslash}p{2.1cm}>{\raggedright\arraybackslash}p{2.1cm}>{\raggedright\arraybackslash}p{2.1cm}}
\toprule
Analysis & Natural history, constant & Natural history, growing & Placebo response & Baseline severity\\
\midrule
Binary-indicator analysis (random intercept, model-based) & \textbf{0.108} & 0.056 & \textbf{0.157} & 0.064\\
Random intercept + CAR(1), CR2 & \textbf{0.103} & 0.049 & \textbf{0.105} & 0.052\\
The same, half-life chosen by AIC & \textbf{0.130} & 0.059 & \textbf{0.125} & 0.057\\
Unstructured covariance, Kenward-Roger & \textbf{0.413} & \textbf{0.200} & \textbf{0.224} & 0.048\\
The same, half-life chosen by AIC & \textbf{0.506} & \textbf{0.258} & \textbf{0.270} & 0.054\\
Random intercept + CAR(1) with expectancy terms, CR2 & 0.045 & 0.040 & 0.052 & 0.055\\
The same, with phase-specific residual variances, CR2 & 0.056 & 0.037 & 0.051 & 0.059\\
The same, with phase-specific residual variances, model-based & 0.054 & 0.042 & 0.049 & 0.059\\
Expectancy terms, half-life chosen by AIC, CR2 & 0.058 & 0.050 & 0.060 & 0.060\\
Random intercept + CAR(1) with free baseline slope, CR2 & 0.051 & 0.066 & 0.060 & 0.054\\
\bottomrule
\end{tabular}
\end{table}

Analyses that keep baseline as an outcome row use one biomarker main
effect for baseline and for the off-drug visits. When the biomarker is
prognostic through a pathway that is absent at baseline, those two
associations differ, and part of the difference is attributed to the
interaction (Appendix \ref{sec-shared}). Every analysis that shares
the biomarker effect with baseline exceeded the nominal level in at
least one such cell (Table \ref{tab:tab-tenan}).

The expectancy-term family stayed within the single-cell two-MCSE band
in every prognostic cell. Its four analyses rejected
0.037 to
0.060. This family adds
the expectancy and a biomarker-by-expectancy term to the mean model,
with or without residual variances that differ by trial phase. The
free-baseline model rejected 0.066
under the growing natural-history pathway, above the single-cell band,
and 0.060 under the placebo pathway.
Appendix \ref{sec-shared} explains why freeing the baseline slope
alone is not sufficient when the prognostic slope grows over time or
scales with the expectancy. A prognostic association with baseline
severity alone did less harm, because baseline severity enters every
visit equally. The binary-indicator analysis rejected
0.064 there. The scenarios in which the
pathway is absent at baseline are the ones that show the inflation.
Whether such a pathway, at a magnitude of 0.3, is realistic for blood
pressure in PTSD is not known.

\textbf{Changes between runs.} Treating two runs as independent, the
difference between them in one cell has a Monte Carlo SE of about
0.010 at a rate of 0.05. The first and third runs use the same seeds
and allocation, so their differences are paired and that SE overstates
their noise. The carryover rule changed the table by at most
0.017, for the expectancy-term model with phase-specific variances (CR2) under the baseline-severity pathway (Appendix \ref{sec-firstrun}). Cells that
crossed the bold threshold of 0.075:
none. Cells that crossed the upper edge
of the single-cell band: the free-baseline model under the growing natural-history pathway (first run 0.063, third run 0.066); the unstructured model with AIC-chosen half-life under the baseline-severity pathway (first run 0.064, third run 0.054).
Randomized allocation changed it by at most 0.041, for the unstructured model with AIC-chosen half-life under the placebo pathway.
Cells that crossed the bold threshold:
the binary-indicator analysis under the growing natural-history pathway (second run 0.077, third run 0.056). Cells that crossed the upper edge
of the single-cell band: the binary-indicator analysis under the growing natural-history pathway (second run 0.077, third run 0.056); the CAR(1) model with AIC-chosen half-life under the growing natural-history pathway (second run 0.068, third run 0.059); the unstructured model with AIC-chosen half-life under the baseline-severity pathway (second run 0.065, third run 0.054).
In all three runs the expectancy-term family stayed within the
single-cell band in every prognostic cell. The binary-indicator, CAR(1)
and unstructured analyses each exceeded it in at least one cell, with
or without an AIC-chosen half-life. Taken analysis by analysis, only
the free-baseline model changed its classification. It stayed within
the band in the first run and exceeded it under the growing pathway in
the other two, by at most
0.009.

\textbf{G10 under these pathways.} The planning sample sizes for G10
(paper 02) assume that the biomarker is not prognostic. We fitted G10 to the
trials of this stress test, with the baseline row. It rejected the
null at 0.052 in the base
cell, and at
0.108,
0.051,
0.110 and
0.050 under the constant natural-history,
growing natural-history, placebo and baseline-severity pathways
(single-cell band 0.036 to 0.064).
The expectancy-term model rejected
0.040 to 0.055 in the
same cells. In this test G10 held its level when the biomarker acted
through baseline severity or the growing natural history, but not when
it acted through a constant natural history or the placebo response.
Section \ref{sec-reconcile} shows why the growing case may not carry
over to an analysis without a baseline row. A trial in which a
constant natural-history or placebo pathway is plausible should be
planned for the expectancy-term analysis. On the power-curve trials of
the stress test, matching G10's power with that analysis needs about
1.13 to 1.27
times the G10 sample size, by the normal approximation to power, at
\(c_{bm}\) from 0.225 to 0.45. At \(c_{bm} = 0.15\) both analyses have
power near 0.15 and the factor,
0.96, is within
Monte Carlo error of 1. Paper 02's G10 sample sizes can be scaled by
that factor. The factor comes from one design, half-life and \(\rho\),
the stress test's base cell.

\textbf{Power and robustness.} The expectancy-term model keeps its size
under every prognostic pathway and has lower power. In the base cell
at \(c_{bm} = 0.30\) its power was
0.404 against
0.486 for the CAR(1) model. A
further 1,500 trials per cell under the second run's DGP (randomized
allocation) were pooled with that run (2,000 per cell). There,
size-adjusted power was 0.431 against
0.496, about
13\%
less. Under compound symmetry at \(c_{bm} = 0.25\) the gap closes. Power
was 0.564 against
0.552 in the reported run, and
size-adjusted power in the pooled trials was
0.542 against
0.539 (Monte Carlo SE about 0.011). The
choice between power and robustness is the trialist's. It depends on
whether a prognostic pathway absent at baseline is plausible for the
biomarker in hand.

\subsubsection{The growing pathway and the trend pathway}\label{sec-reconcile}

The growing natural-history pathway of this test and the trend
pathway of Section \ref{sec-trend} are the same mean shift. Both add
\(0.3\,\sigma_{TV}\,z_i\,g(w_t)\) to the natural-history component,
with \(g\) the natural-history mean course scaled to peak 1 over the
visits. Driver 11 applies it in \texttt{draw()}. The paper 02 driver
applies it in \texttt{implementations/tidyverse/R/functions.R}. Both use the
published natural-history curve, \(\sigma_{TV} =\) 10,
the Hybrid design with fixed allocation, \(N = 70\), \(t_{1/2} = 0.5\)
week and 1,000 null replicates. The two results appear to disagree.
Under the trend pathway G9 exceeded the band. Under the growing
pathway G10 stayed inside it.

The comparison is less direct than that. The two tests did not fit
the same analyses. Block T fitted G8, G9 and F1, not G10. The second
test fitted G10 but not G8 or G9. Its CAR(1) model with CR2 at the
fixed half-life of 0.5 week has the formula, covariance and test of
G8. Its AIC variant resembles G9, with a grid that adds 0. The two
tests also differ in their analysis data. Every analysis of the
second test has a baseline row. The paper 02 driver has none. Table
\ref{tab:tab-reconcile} sets the comparable analyses side by side.
The comparison is across jobs and unpaired.

\begin{table}
\centering
\caption{\label{tab:tab-reconcile}The growing natural-history shift ($0.3\,
      \sigma_{TV}\,z_i\,g(w_t)$) with no true interaction, in the two
      drivers. Each entry: rejection rate (mean interaction estimate
      per SD of the biomarker; empirical SE). Hybrid, $N = 70$,
      $t_{1/2} = 0.5$ week, fixed allocation, 1,000 replicates. Band
      $[0.036, 0.064]$. The estimates of the second test are converted from raw
      biomarker units by its SD.}
\centering
\fontsize{8}{10}\selectfont
\begin{tabular}[t]{>{\raggedright\arraybackslash}p{4cm}>{\raggedright\arraybackslash}p{4.6cm}>{\raggedright\arraybackslash}p{4.6cm}}
\toprule
Analysis & Block T, no baseline row & Second stress test, baseline row\\
\midrule
Fixed half-life 0.5 week, CR2 & G8: 0.063 (+0.43; 1.45) & CAR(1): 0.049 (-0.14; 1.22)\\
Half-life by AIC, interaction in the model & G9: 0.073 (+0.42; 1.46) & AIC variant: 0.059 (-0.19; 1.28)\\
Half-life by AIC, interaction not in the model & not fitted & G10: 0.051 (-0.18; 1.24)\\
Expectancy terms, fixed half-life & F1 (0.2 week): 0.052 (+0.24; 1.27) & Expectancy-term (0.5 week): 0.040 (+0.03; 1.37)\\
\bottomrule
\end{tabular}
\end{table}

At the fixed half-life the two drivers agree on size. G8 rejected at
0.063 and the CAR(1) model at 0.049, both inside
the band. G9 rejected at 0.073, above it, and the AIC
variant of the second test at 0.059. The mean estimates
differ in sign. Without a baseline row G8 gave +0.43 per SD
of the biomarker. With one, the CAR(1) model gave -0.14 and
G10 -0.18. Every one of these is within a third of its
empirical SE of zero.

The baseline row accounts for the change of sign. The shift is zero
at week 0 and grows along the natural-history curve. In the Hybrid
design the off-drug visits come late, where the shift is large.
Without a baseline row the off-drug visits carry more of the shift
than the on-drug visits. The shift lowers the outcome, so the
difference projects onto \(B^c_i D_{bc,it}\) with a positive sign. A
baseline row adds an off-drug row with no shift. That reverses the
sign. An unweighted least-squares projection of the shift onto the
biomarker terms of the G8 model gives +0.69 per SD of the
biomarker without the baseline row and -0.08 with it. It
ignores the random intercept and the CAR(1) weighting, so it checks
only the signs, which agree with the fitted estimates.

The two results therefore do not conflict. The DGP settings match.
The difference lies in the analysis data and in which analyses were
fitted. Two points remain open. First, G10 has not been fitted under
the trend pathway without a baseline row. That it held its level
under the growing pathway rests on analyses with a baseline row only.
Second, the bias is real in both tests. Every G-type mean estimate in
Table \ref{tab:tab-reconcile} is at least 3.6 Monte
Carlo SEs from zero. It is within a third of the empirical SE at
\(N = 70\), so it moves the rejection rate little there. As under the
placebo-response pathway (Section \ref{sec-largen}), a fixed bias of
either sign inflates the size of the two-sided test as \(N\) grows. The
sign here depends on the visit structure of the analysis. A larger
sample under the trend pathway, with the same analysis fitted with
and without the baseline row, would settle both points. The biomarker-by-time block of Section \ref{sec-bmt}
should include that contrast.

\subsection{The repair artifact: a historical exhibit}\label{sec-artifact}

\fbox{\parbox{0.95\linewidth}{\textbf{Historical exhibit.} Every
number in this section comes from runs under the simcore2026 parameter
set with silent repair (Table \ref{tab:tab-sources}). They are
evidence about a repaired DGP. They are not evidence about any
analysis under a prognostic biomarker.}}

The first prognostic stress test ran at strengths 0.3 and 0.6 before
the drivers refused infeasible matrices. At strength 0.6 on the
placebo-response and natural-history pathways, G9 lost
12 to 26 points of power at \(c_{bm} = 0.45\) against the
no-confound reference of 0.759 (\(N = 35\)) and
0.969 (\(N = 70\)). At strength 0.3 the largest loss was
0.039. Every one of the 14 analyses fitted
to the same trials lost power at strength 0.6 (Table
\ref{tab:tab-hist}). That uniformity was the first sign that the
loss arose before the analysis.

\begin{table}
\centering
\caption{\label{tab:tab-hist}Historical exhibit (simcore2026 parameter set,
      silent repair): power at $c_{bm} = 0.45$ at prognostic
      strengths 0.3 and 0.6, Hybrid, $t_{1/2} = 0.5$ week, 500
      replicates. Every strength-0.6 cell here, including
      $c_{bm,BL} = 0.6$, was simulated under a repaired matrix.}
\centering
\fontsize{9}{11}\selectfont
\begin{tabular}[t]{lrllllll}
\toprule
\multicolumn{2}{c}{ } & \multicolumn{2}{c}{G9} & \multicolumn{2}{c}{F1} & \multicolumn{2}{c}{F1A0B} \\
\cmidrule(l{3pt}r{3pt}){3-4} \cmidrule(l{3pt}r{3pt}){5-6} \cmidrule(l{3pt}r{3pt}){7-8}
Pathway & $N$ & 0.3 & 0.6 & 0.3 & 0.6 & 0.3 & 0.6\\
\midrule
$c_{bm,PB}$ & 35 & 0.798 & 0.597 & 0.728 & 0.524 & 0.741 & 0.562\\
 & 70 & 0.968 & 0.853 & 0.946 & 0.798 & 0.940 & 0.794\\
$c_{bm,TV}$ & 35 & 0.782 & 0.502 & 0.720 & 0.454 & 0.730 & 0.482\\
 & 70 & 0.968 & 0.772 & 0.952 & 0.718 & 0.950 & 0.702\\
$c_{bm,BL}$ & 35 & 0.782 & 0.738 & 0.713 & 0.676 & 0.758 & 0.702\\
 & 70 & 0.976 & 0.950 & 0.958 & 0.929 & 0.960 & 0.920\\
\bottomrule
\end{tabular}
\end{table}

A direct check of the covariance matrices explains it (Table
\ref{tab:tab-pd}). Each strength-0.6 cell on those two pathways, and
\(c_{bm,BL} = 0.6\) at \(c_{bm} = 0.45\), has a joint matrix that is not
positive definite on all 4 Hybrid paths. The repair
replaced it with a nearby matrix before any data were drawn. The
repair shrank the population on-drug minus off-drug slope of the
outcome on the biomarker. That is the quantity the interaction term
estimates. With no prognostic correlation it is -3.26.
After repair it was -2.38 under the placebo-response
pathway and -1.81 under the natural-history pathway. The
power loss is the repair removing part of the interaction from the
data. No analysis could recover it. At strength 0.3 the matrix is
positive definite, with a smallest eigenvalue of 0.010
under the placebo-response pathway.

\begin{table}
\centering
\caption{\label{tab:tab-pd}Historical exhibit: positive definiteness of the raw
      joint covariance matrix on the Hybrid paths, and the population
      on-drug minus off-drug slope of the outcome on the biomarker
      before and after the silent repair. Computed from the matrix
      under the simcore2026 SDs, no simulation
      (\texttt{30-check-prognostic-confound-pd.R}).}
\centering
\fontsize{9}{11}\selectfont
\begin{tabular}[t]{lllrlll}
\toprule
Pathway & Strength & $c_{bm}$ & Paths repaired & Min. eigenvalue & Contrast, raw & Contrast, repaired\\
\midrule
none & 0.0 & 0.45 & 0 & 0.180 & -3.26 & -3.26\\
$c_{bm,BL}$ & 0.3 & 0.45 & 0 & 0.145 & -3.26 & -3.26\\
$c_{bm,BL}$ & 0.6 & 0.00 & 0 & 0.184 & +0.00 & +0.00\\
$c_{bm,BL}$ & 0.6 & 0.45 & 4 & -0.110 & -3.26 & -3.00\\
$c_{bm,PB}$ & 0.3 & 0.45 & 0 & 0.010 & -3.40 & -3.40\\
$c_{bm,PB}$ & 0.6 & 0.00 & 4 & -0.240 & -0.27 & -0.23\\
$c_{bm,PB}$ & 0.6 & 0.45 & 4 & -0.642 & -3.53 & -2.38\\
$c_{bm,TV}$ & 0.3 & 0.45 & 0 & 0.011 & -3.26 & -3.26\\
$c_{bm,TV}$ & 0.6 & 0.00 & 4 & -0.354 & +0.00 & +0.00\\
$c_{bm,TV}$ & 0.6 & 0.45 & 4 & -0.853 & -3.26 & -1.81\\
\bottomrule
\end{tabular}
\end{table}

Before the artifact was found, the loss was read as patient-level
variance that a random intercept could not absorb. G9 was refitted
with a random slope on time, on expectancy, and on both, in
12 cells at 1,000 replicates. 8
of those cells are repaired strength-0.6 cells. The three variants
matched plain G9 to within 0.010 in rejection rate in every
cell (Table \ref{tab:tab-re}). The largest paired per-replicate
difference was 0.008. Non-convergence reached
2.5 percent for the variant with both slopes. The DGP
has no patient-specific slope, so there was no variance for a random
slope to absorb. The check shows only that the slopes cost nothing
there.

\begin{table}
\centering
\caption{\label{tab:tab-re}Historical exhibit: rejection rate of G9 and
      its three random-slope variants (simcore2026 parameter set,
      Hybrid, $t_{1/2} = 0.5$ week, 1,000 replicates). The strength-0.6
      rows were simulated under repaired matrices.}
\centering
\fontsize{9}{11}\selectfont
\begin{tabular}[t]{lllll}
\toprule
Cell & G9 & + slope on $w_t$ & + slope on $\eta_t$ & + both\\
\midrule
none, $c_{bm} = 0.00$, $N = 35$ & 0.050 & 0.050 & 0.050 & 0.046\\
none, $c_{bm} = 0.00$, $N = 70$ & 0.043 & 0.044 & 0.043 & 0.041\\
none, $c_{bm} = 0.45$, $N = 35$ & 0.759 & 0.762 & 0.766 & 0.769\\
none, $c_{bm} = 0.45$, $N = 70$ & 0.969 & 0.970 & 0.969 & 0.970\\
$c_{bm,PB}$ 0.6, $c_{bm} = 0.00$, $N = 35$ & 0.060 & 0.060 & 0.059 & 0.060\\
$c_{bm,PB}$ 0.6, $c_{bm} = 0.00$, $N = 70$ & 0.047 & 0.046 & 0.047 & 0.047\\
$c_{bm,PB}$ 0.6, $c_{bm} = 0.45$, $N = 35$ & 0.594 & 0.593 & 0.593 & 0.596\\
$c_{bm,PB}$ 0.6, $c_{bm} = 0.45$, $N = 70$ & 0.860 & 0.859 & 0.858 & 0.858\\
$c_{bm,TV}$ 0.6, $c_{bm} = 0.00$, $N = 35$ & 0.050 & 0.053 & 0.051 & 0.054\\
$c_{bm,TV}$ 0.6, $c_{bm} = 0.00$, $N = 70$ & 0.044 & 0.044 & 0.043 & 0.040\\
$c_{bm,TV}$ 0.6, $c_{bm} = 0.45$, $N = 35$ & 0.489 & 0.488 & 0.489 & 0.485\\
$c_{bm,TV}$ 0.6, $c_{bm} = 0.45$, $N = 70$ & 0.772 & 0.772 & 0.773 & 0.768\\
\bottomrule
\end{tabular}
\end{table}

The published-set rerun applied the current default, which refuses an
infeasible matrix. Of its 12 strength-0.6 cells on the
three correlation pathways, 10 were refused and
2 ran. The refused cells are exactly the configurations
that the historical check found infeasible. The cells that ran set
\(c_{bm,BL} = 0.6\) with no true interaction, which is feasible. Under
the current code a repaired cell cannot be simulated without an
explicit option.

\section{Discussion}\label{sec-discussion}

\subsection{What the expectancy terms are for}\label{sec-expectancy}

The placebo-response amplitude scales with \(\eta_t\). In the Hybrid
design \(\eta_t\) differs between on-drug and off-drug visits. A
biomarker correlated with the placebo response therefore acquires an
on/off slope difference with no true interaction. Under the
simcore2026 SDs the population contrast was -0.13 at
\(c_{bm,PB} = 0.3\) and -0.27 at 0.6, against
-3.26 for the interaction at \(c_{bm} = 0.45\) (historical,
Table \ref{tab:tab-pd}). The placebo-response SD is
2 in the simcore2026 set and 10 in the
published set. The contrast scales with that SD. It has not been
recomputed under the published set.

The published-set runs show the consequence. Under the
placebo-response pathway the G-family estimate moved away from zero
with no interaction present (Section \ref{sec-grid}). The bias was
small against the standard error at \(N = 70\) and produced no excess
rejection there. At \(N = 280\) it did (Section \ref{sec-largen}).
The F-family, whose \(B^c_i\eta_t\) term absorbs the contrast, stayed
inside the band at every \(N\) examined. The earlier conclusion that
the expectancy terms add nothing came from the simcore2026 runs. It
does not hold under the published parameter values.

\subsection{The joint recommendation}\label{sec-recommend}

The two companion papers give one recommendation between them.

\begin{itemize}
\tightlist
\item
  For a biomarker believed to be predictive only, use G10: the
  exposure-weighted indicator with the half-life chosen by AIC on the
  model without the interaction, tested with CR2
  (\citeproc{ref-pmsimstats-paper02}{3}).
\item
  Where the biomarker may also be prognostic, add the
  design-expectancy terms \(\eta_t\) and \(B^c_i\eta_t\). This is F1A0,
  or its bootstrap-calibrated form F1A0B. F1A0 and G10 differ only in
  those two terms.
\item
  Where the biomarker may predict the rate of natural-history change,
  also add a biomarker-by-time term. This step is provisional until
  the block of Section \ref{sec-bmt} is run.
\end{itemize}

F1, which fixes the half-life at 0.2 week, is not recommended. It
loses power wherever that value is wrong (Section \ref{sec-gvf}).
Two questions remain for F1A0B. Its bootstrap calibration may be
unnecessary once the selection excludes the interaction. Its cost
against G10 without a prognostic effect is measured here on one grid
only.

\subsection{Limitations}\label{sec-limitations}

The published-set results come from one parameter set and one
correlation structure (\(\rho = 0.7\), \(c_1 = 0.2\)). Most cells are on
the Hybrid design at \(t_{1/2} = 0.5\) week. The prognostic correlations
are correlations of a single latent biomarker with the components.
The feasibility bound is a property of that parameterization. The
large-sample cells have no confound-free counterpart at the same \(N\).
The trend pathway was simulated at one coefficient, one sample size
and one half-life. The sign of its bias depends on whether the
analysis has a baseline row, and no single job has fitted the same
analysis both ways. G10 has not been fitted under it without a
baseline row. Its remedy has not been simulated. The second stress
test ran 1,000 null replicates per cell at one sample size, and its
cells and analyses were revised as results came in. The
placebo-response contrast has not been recomputed under the published
SDs. The random-slope variants were fitted only to historical,
partly repaired cells under a DGP without patient-specific slopes.
Comparisons with the block A reference are across jobs and unpaired.
Cells within a job share random-number streams, so counts of cells
inside a band are not independent confirmations.

\section{Conclusions}\label{sec-conclusions}

\begin{enumerate}
\def\labelenumi{\arabic{enumi}.}
\tightlist
\item
  A constant prognostic shift is absorbed exactly by the biomarker
  main effect. It cannot move the interaction estimate.
\item
  Inside the feasibility bounds, the correlation pathways left the G9
  and G10 Type I error inside the band at \(N \le 70\).
\item
  A biomarker correlated with the placebo response biases the
  G-family estimate. Under the published parameter values the bias
  inflated the size of G9 and G10 at \(N = 280\). The expectancy terms
  removed it.
\item
  A biomarker that predicts the natural-history trend biased every
  analysis fitted by the paper 02 driver, which has no baseline row.
  The expectancy terms absorbed part of the bias. With a baseline
  row the bias changed sign. At \(N = 70\) the mean estimate was
  within a third of its empirical SE of zero in both. The
  biomarker-by-time remedy is pending.
\item
  In the second stress test, whose analyses have a baseline row, the
  expectancy-term analyses kept their size under all four pathways.
  G10 exceeded the band under the constant natural-history and
  placebo-response pathways.
\item
  The power loss once attributed to a prognostic correlation of 0.6
  was produced by the silent repair of infeasible matrices. Under the
  current default those cells are refused.
\item
  For a biomarker that may be prognostic, add \(\eta_t\) and
  \(B^c_i\eta_t\) to G10. The result is F1A0.
\end{enumerate}

\section*{Declarations}\label{declarations}
\addcontentsline{toc}{section}{Declarations}

\textbf{Funding, competing interests, author contributions.} To be
completed.

\textbf{Data and code availability.} Every number in this draft is read
from the stored files of Table \ref{tab:tab-sources} by inline code.

\section{Reproducibility}\label{sec-repro}

The drivers live in \texttt{analysis/scripts/carryover-sensitivity/}. The
published-set runs were submitted on 2026-11-03 under \texttt{docs/74},
wave 1. The prognostic stress test is driver 28. The feasible grid and
the large-sample cells are driver 32, block C. The trend, null and
alternative blocks are driver 36, blocks T, N and A. The original
grid is driver 27 and the scenario cells driver 06. Driver 29, the
random-slope variants, was not rerun. Every confounded cell in it is
a repaired strength-0.6 cell (\texttt{docs/74}, decision D5).

The second stress test is study 2 of paper 01. Its driver is
\texttt{analysis/scripts/quick-sim/carryover-closed-form/11-strawman-stress-test.R},
and its outputs are under
\texttt{analysis/data/quick-sim/carryover-closed-form/strawman-stress/}. The
driver holds the published means and SDs as constants and stops on a
matrix that is not positive definite. The reported (third) run is
reproduced with \texttt{-\/-carryover\ anchored\ -\/-allocation\ fixed\ -\/-analysis-half-life\ 0.5} (files tagged \texttt{-anchored-fixed-ha0.5}). The
second run used \texttt{-\/-carryover\ anchored\ -\/-allocation\ random\ -\/-analysis-half-life\ 0.5} (tagged \texttt{-anchored}). The first used
\texttt{-\/-carryover\ recursive\ -\/-allocation\ fixed\ -\/-analysis-half-life\ 0.5}
(untagged). Each run was followed by the same command with \texttt{-\/-extra}
and with \texttt{-\/-extra2} added, on the same trials. G10 was fitted to the
trials of the third run by the same command with \texttt{-\/-extra5}. The
pooled size-adjusted comparison adds 1,500 fresh trials per cell from
\texttt{-\/-extra4} under the second run's DGP (tagged \texttt{-extra4-anchored});
\texttt{15-extra4-pooled-inputs.R} pools them with the second run and writes
\texttt{strawman-stress/extra4-decision/pooled/}. The script's default analysis half-life is now
0.2 week, so the \texttt{-\/-analysis-half-life\ 0.5} flag is needed. The joint
feasibility bounds of Table \ref{tab:tab-jointceil} come from
\texttt{analysis/scripts/quick-sim/carryover-closed-form/18-prognostic-joint-ceiling.R},
which writes
\texttt{analysis/data/quick-sim/carryover-closed-form/rho-stage/prognostic-joint-ceilings.csv}.
It computes the bounds from the correlation matrix, with no
simulation.

Changes needed before the next draft, not made here:

\begin{itemize}
\tightlist
\item
  Add the biomarker-by-time specifications (\texttt{E2cr2\_bmt},
  \texttt{E7cr2\_bmt}, \texttt{E7null\_cr2\_bmt}, \texttt{F1A0\_bmt}, optionally quadratic) to
  \texttt{simulation-core.R}. Add the corresponding block to
  \texttt{36-run-referee-variants.R} (split plan, Section 3.1, step 6).
\item
  Give \texttt{30-check-prognostic-confound-pd.R} a \texttt{-\/-params} switch so that
  the placebo-response contrast can be recomputed under the published
  SDs.
\item
  Give \texttt{37-ladder-tables-figures.R\ -\/-scope\ all} a published-set input
  path if paper 15 is to carry ladder-style figures. Its trend figure
  (\texttt{02-ladder-fig4-trend.png}) was built from historical data and is
  not used here.
\end{itemize}

\section*{References}\label{references}
\addcontentsline{toc}{section}{References}

\protect\phantomsection\label{refs}
\begin{CSLReferences}{0}{1}
\bibitem[\citeproctext]{ref-hendrickson2020}
\CSLLeftMargin{1. }%
\CSLRightInline{Hendrickson RC, Thomas RG, Schork NJ, Raskind MA. Optimizing aggregated {N}-of-1 trial designs for predictive biomarker validation: Statistical methods and theoretical findings. Frontiers in Digital Health. 2020;2:13. doi:\href{https://doi.org/10.3389/fdgth.2020.00013}{10.3389/fdgth.2020.00013}}

\bibitem[\citeproctext]{ref-pmsimstats-paper01}
\CSLLeftMargin{2. }%
\CSLRightInline{{pmsimstats team}. Simulating predictive-biomarker interactions in aggregated {N}-of-1 trials: An extended data-generating process building on {Hendrickson} et al. 2026.}

\bibitem[\citeproctext]{ref-pmsimstats-paper02}
\CSLLeftMargin{3. }%
\CSLRightInline{{pmsimstats team}. Carryover representation and inference for biomarker-treatment interaction tests with a predictive biomarker in aggregated {N}-of-1 trials under covariance moderation: A staged simulation study and a recommended analysis for the {Hybrid} design. 2026.}

\bibitem[\citeproctext]{ref-bell2002bias}
\CSLLeftMargin{4. }%
\CSLRightInline{Bell RM, McCaffrey DF. Bias reduction in standard errors for linear regression with multi-stage samples. Survey Methodology. 2002;28(2):169--81.}

\bibitem[\citeproctext]{ref-pustejovsky2018}
\CSLLeftMargin{5. }%
\CSLRightInline{Pustejovsky JE, Tipton E. Small-sample methods for cluster-robust variance estimation and hypothesis testing in fixed effects models. Journal of Business and Economic Statistics. 2018;36(4):672--83. doi:\href{https://doi.org/10.1080/07350015.2016.1247004}{10.1080/07350015.2016.1247004}}

\bibitem[\citeproctext]{ref-brookes2004}
\CSLLeftMargin{6. }%
\CSLRightInline{Brookes ST, Whitely E, Egger M, Davey Smith G, Mulheran PA, Peters TJ. Subgroup analyses in randomized trials: Risks of subgroup-specific analyses; power and sample size for the interaction test. Journal of Clinical Epidemiology. 2004;57(3):229--36. doi:\href{https://doi.org/10.1016/j.jclinepi.2003.08.009}{10.1016/j.jclinepi.2003.08.009}}

\bibitem[\citeproctext]{ref-senn2002}
\CSLLeftMargin{7. }%
\CSLRightInline{Senn S. Cross-over trials in clinical research. 2nd ed. Chichester, UK: John Wiley \& Sons; 2002. doi:\href{https://doi.org/10.1002/0470854596}{10.1002/0470854596}}

\bibitem[\citeproctext]{ref-araujo2016}
\CSLLeftMargin{8. }%
\CSLRightInline{Araujo A, Julious S, Senn S. Understanding variation in sets of {N}-of-1 trials. PLoS ONE. 2016;11(12):e0167167. doi:\href{https://doi.org/10.1371/journal.pone.0167167}{10.1371/journal.pone.0167167}}

\bibitem[\citeproctext]{ref-haller2018}
\CSLLeftMargin{9. }%
\CSLRightInline{Haller B, Ulm K. A simulation study on estimating biomarker-treatment interaction effects in randomized trials with prognostic variables. Trials. 2018;19(1):128.}

\bibitem[\citeproctext]{ref-kent2020path}
\CSLLeftMargin{10. }%
\CSLRightInline{{Kent DM, Paulus JK, Klaveren D van, et al.} The predictive approaches to treatment effect heterogeneity ({PATH}) statement. Annals of Internal Medicine. 2020;172(1):35--45. doi:\href{https://doi.org/10.7326/M18-3667}{10.7326/M18-3667}}

\bibitem[\citeproctext]{ref-akaike1974new}
\CSLLeftMargin{11. }%
\CSLRightInline{Akaike H. A new look at the statistical model identification. IEEE Transactions on Automatic Control. 1974;19(6):716--23. doi:\href{https://doi.org/10.1109/TAC.1974.1100705}{10.1109/TAC.1974.1100705}}

\bibitem[\citeproctext]{ref-berk2013valid}
\CSLLeftMargin{12. }%
\CSLRightInline{Berk R, Brown L, Buja A, Zhang K, Zhao L. Valid post-selection inference. The Annals of Statistics. 2013;41(2):802--37. doi:\href{https://doi.org/10.1214/12-AOS1077}{10.1214/12-AOS1077}}

\bibitem[\citeproctext]{ref-leeb2005}
\CSLLeftMargin{13. }%
\CSLRightInline{Leeb H, Pötscher BM. Model selection and inference: Facts and fiction. Econometric Theory. 2005;21(1):21--59. doi:\href{https://doi.org/10.1017/S0266466605050036}{10.1017/S0266466605050036}}

\bibitem[\citeproctext]{ref-morris2019simulation}
\CSLLeftMargin{14. }%
\CSLRightInline{Morris TP, White IR, Crowther MJ. Using simulation studies to evaluate statistical methods. Statistics in Medicine. 2019;38(11):2074--102. doi:\href{https://doi.org/10.1002/sim.8086}{10.1002/sim.8086}}

\bibitem[\citeproctext]{ref-liu2009baseline}
\CSLLeftMargin{15. }%
\CSLRightInline{Liu GF, Lu K, Mogg R, Mallick M, Mehrotra DV. Should baseline be a covariate or dependent variable in analyses of change from baseline in clinical trials? Statistics in Medicine. 2009;28(20):2509--30. doi:\href{https://doi.org/10.1002/sim.3639}{10.1002/sim.3639}}

\end{CSLReferences}

\appendix


\section{Absorption of a constant shift}\label{sec-absorb}

Let \(X\) be the fixed-effects design matrix of any analysis here. It
contains the biomarker main effect \(B^c_i\). The constant shift adds
\(b_{TV}\sigma_{TV} z_i\) to every outcome of participant \(i\). Since
\(z_i\) is an affine function of \(B^c_i\), the shift vector lies in the
column space of \(X\). For fixed variance parameters the generalized
least-squares estimate then changes only in the coefficients of the
intercept and \(B^c_i\). The residuals are unchanged. REML and ML
estimates of the variance parameters depend on the data only through
those residuals, so they are unchanged too. AIC selection compares
fits whose likelihoods are unchanged, so it selects the same
half-life. CR2 standard errors depend on the residuals and the design
only. The interaction estimate, its standard error and its \(p\)-value
are therefore identical, replicate by replicate. The stored data
confirm it (Section \ref{sec-stress}).

\section{The joint feasibility condition}\label{sec-jointfeas}

Order the variables as biomarker, then everything else, and let \(M_0\)
be the correlation matrix of everything but the biomarker, including
the baseline level, which is uncorrelated with the response
components. By the Schur complement the joint matrix is positive
definite if and only if \(M_0\) is and the biomarker row \(r\) satisfies
\(r^\top M_0^{-1} r < 1\). With \(c_{bm}\) alone, \(r = c_{bm}\tilde g\),
where \(\tilde g\) is the correlation profile padded with zeros, and the
bound is \(c_{bm}^{\ast} = (\tilde g^\top M_0^{-1}\tilde g)^{-1/2}\)
(\citeproc{ref-pmsimstats-paper01}{2}).

With prognostic correlations the row is a sum of fixed patterns scaled
by the correlations,
\(r = c_{bm}\tilde g + c_{bm,PB}\, p + c_{bm,TV}\, q + c_{bm,BL}\, b\).
Here \(p\) and \(q\) are one on every visit of the placebo-response and
natural-history components, and \(b\) is one on the baseline level. The
condition becomes \(c^\top A c < 1\), with
\(c = (c_{bm}, c_{bm,PB}, c_{bm,TV}, c_{bm,BL})\) and
\(A_{jk} = x_j^\top M_0^{-1} x_k\) for \(x \in \{\tilde g, p, q, b\}\). The
feasible region is the interior of an ellipsoid. Fixing all
correlations but one makes \(c^\top A c = 1\) a quadratic in that
correlation, whose positive root is its bound. A negative discriminant
means the fixed correlations are already infeasible. Because \(b\) is an
eigenvector of \(M_0\) with unit eigenvalue, \(c_{bm,BL}\) enters only
through its square. Its bound given \(c_{bm}\) is
\(\sqrt{1 - c_{bm}^2 A_{\tilde g\tilde g}}\). The \(p\) and \(q\) patterns
are structurally identical, so \(c_{bm,PB}\) and \(c_{bm,TV}\) have the
same bounds. A design's bound is the minimum over its paths. Table
\ref{tab:tab-jointceil} gives the values for the Hybrid design. The
off-diagonal entries of \(A\) are negative at the extended DGP's values,
so the correlations are mildly complementary. The DGP sets the
biomarker's marginal correlations with the other components to zero.
The negative entries of \(A\) reflect partial association: given its
correlation with the drug response, the biomarker's partial
correlation with the other components is not zero. An explicit
correlation of the same sign therefore uses less of the feasible
region than it would if the components were independent.

\section{The shared biomarker effect}\label{sec-shared}

Let \(s_0\), \(s_{\text{off}}\) and \(s_{\text{on}}\) be the true slopes of
\(Y\) on \(B\) at baseline, at post-baseline off-drug visits and at
on-drug visits. A model with baseline as an off-drug row estimates one
\(\beta_B\) for \(s_0\) and \(s_{\text{off}}\).

The following account is heuristic, not a derivation. Consider a
prognostic biomarker with no interaction, whose post-baseline slope is
constant: \(s_0 = 0 \ne s_{\text{off}} = s_{\text{on}} = s_{\text{prog}}\).

\begin{itemize}
\tightlist
\item
  \textbf{The bias.} The fitted \(\hat\beta_B\) lies between the baseline and
  off-drug slopes, \(\hat\beta_B \approx f\,s_{\text{prog}}\) for some
  \(0 < f < 1\) that depends on the weight the baseline row receives.
  The interaction absorbs the remainder,
  \(\hat\beta_{bm:D} \approx (1 - f)\,s_{\text{prog}} \ne 0\).
\item
  \textbf{The correction.} Adding \(B_i\cdot\mathbf 1\{t = 0\}\) frees \(s_0\)
  and removes this source of bias. The constraint is analogous to the
  choice, in longitudinal analysis, between treating baseline as a
  covariate and as a response (\citeproc{ref-liu2009baseline}{15}).
\end{itemize}

Two further mechanisms operate in the scenarios of Section
\ref{sec-tenan}.

\begin{itemize}
\tightlist
\item
  \textbf{The slope grows over time.} When the post-baseline slope grows,
  the bias depends on how drug status is confounded with time.
  Freeing \(s_0\) is then not sufficient in general. The rejection rate
  of 0.066 for that analysis under the
  growing pathway is consistent with this. Section
  \ref{sec-reconcile} shows that the same growing shift biases the
  interaction with or without a baseline row, in opposite directions.
\item
  \textbf{Expectancy.} When the biomarker is associated with the placebo
  response, the slope also scales with the expectancy \(\eta_t\), which
  is higher at the open-label (on-drug) visits. The expectancy term
  \(B_i \cdot \eta_t\), with \(\eta_0 = 0\), addresses both this mechanism
  and the baseline constraint. Freeing \(s_0\) alone does not. An
  expectancy-scaled slope still enters the interaction through the
  open-label visits.
\end{itemize}

Neither correction is exact in the Hybrid design. The expectancy term
removes a constant post-baseline slope from the interaction only when
the blinded on-drug and off-drug visits carry equal total weight. In
the Hybrid design they do not, so a small bias remains. The
corresponding derivations are in the project's compendium of
mathematical results (\texttt{analysis/report/compendium-mathematics.tex}).
The null rejection rates of Section \ref{sec-tenan} across the four
prognostic cells were 0.040 to
0.055 for the expectancy-term model and
0.051 to 0.066 for the
free-baseline model. The residual bias of the expectancy-term model is
therefore within Monte Carlo error at the prognostic correlations
simulated. That of the free-baseline model is not, under the growing
pathway. Under a prognostic association with baseline severity alone
the binary-indicator analysis rejected
0.064, and the two corrections rejected
0.055 and 0.054.

\section{First run of the second stress test}\label{sec-firstrun}

The first run of the stress test of Section \ref{sec-tenan} used the
compounding carryover rule on the mean and fixed path allocation (18,
18, 17, 17). Its table is given here for comparison with Table
\ref{tab:tab-tenan}. The two runs differ only in the carryover rule.

\begin{table}
\centering
\caption{\label{tab:tab-tenan-first}Second stress test, first run (compounding carryover
    rule, fixed allocation): null rejection rates under a prognostic
    biomarker, 1,000 replicates per cell (two-MCSE band 0.036 to
    0.064 for a single cell). Rates above 0.075 in bold.}
\centering
\fontsize{9}{11}\selectfont
\begin{tabular}[t]{>{\raggedright\arraybackslash}p{5.5cm}>{\raggedright\arraybackslash}p{2.1cm}>{\raggedright\arraybackslash}p{2.1cm}>{\raggedright\arraybackslash}p{2.1cm}>{\raggedright\arraybackslash}p{2.1cm}}
\toprule
Analysis & Natural history, constant & Natural history, growing & Placebo response & Baseline severity\\
\midrule
Binary-indicator analysis (random intercept, model-based) & \textbf{0.107} & 0.053 & \textbf{0.159} & 0.075\\
Random intercept + CAR(1), CR2 & \textbf{0.099} & 0.049 & \textbf{0.100} & 0.050\\
The same, half-life chosen by AIC & \textbf{0.115} & 0.058 & \textbf{0.112} & 0.055\\
Unstructured covariance, Kenward-Roger & \textbf{0.416} & \textbf{0.195} & \textbf{0.222} & 0.051\\
The same, half-life chosen by AIC & \textbf{0.509} & \textbf{0.245} & \textbf{0.256} & 0.064\\
Random intercept + CAR(1) with expectancy terms, CR2 & 0.047 & 0.046 & 0.040 & 0.040\\
The same, with phase-specific residual variances, CR2 & 0.055 & 0.044 & 0.042 & 0.042\\
The same, with phase-specific residual variances, model-based & 0.062 & 0.043 & 0.045 & 0.046\\
Expectancy terms, half-life chosen by AIC, CR2 & 0.059 & 0.052 & 0.047 & 0.045\\
Random intercept + CAR(1) with free baseline slope, CR2 & 0.049 & 0.063 & 0.049 & 0.050\\
\bottomrule
\end{tabular}
\end{table}

\section{Draft consistency checks}\label{sec-checks}

The prose of this draft makes 40 qualitative statements
that are checked against the stored data when it renders. All of them hold.

\end{document}
